\section{Complete results: the Phase 5.5 orthogonalized variants}
\label{sec:complete55}

Section~\ref{sec:complete} reports the whole measurement surface for the
proposed model and its comparators on the frozen tournament coordinates. This
section does the same job for the three orthogonalized variants introduced in
Phase 5.5, which were screened on their own manifest and therefore cannot be
appended as extra columns to the tables above without silently changing what a
column means. Every number here is generated from the audited result files by
\texttt{report/build\_phase55\_assets.py} for the Stage 1 screen and by
\texttt{report/build\_phase55\_stage2\_assets.py} for Stage 2, both of which
read and aggregate but never fit a model, exactly as their companion does.

The three variants replace one component of the estimator and leave the rest of
the pipeline frozen. \texttt{cwdb\_rmean} is a vector R-learner in the sense of
Nie and Wager: it fits the causal contrast by minimising the R-loss
$n^{-1}\sum_i \lVert Z_i - \hat m_{-i}(X_i) - (A_i - \hat e_{-i}(X_i))\,t(X_i)
\rVert^2$ over the rescaled quantile coordinates, with exact weighted
least-squares leaves that never divide by $A - \hat e$. \texttt{cwdb\_mutau} is
the Bayesian-causal-forest reparameterisation of the arm-shared tree, carrying a
prognostic field and a contrast field in each leaf through $z_{a,m}(x) = b_m(x)
+ (a - \hat e(x))\,d_m(x)$. \texttt{cwdb\_xmean} is a vector X-learner with
out-of-fold imputed pseudo-outcomes. All three draw their propensity and
prognostic nuisances from one cross-fitted stack with a declared clip of
$[0.02, 0.98]$ and a canonical fold plan that is a pure function of the data
rather than of row order.

The contract difference between them is not cosmetic and it governs which
tables a variant may appear in. \texttt{cwdb\_mutau} produces a conditional law
and is therefore scored on every metric in the report. \texttt{cwdb\_rmean} and
\texttt{cwdb\_xmean} target the arm contrast only and produce no law at all, so
their entries in the law tables are \emph{n/a} by construction rather than by
omission, and no reading of those tables should treat a missing law column as a
loss.

\paragraph{Conventions, and where they differ from Section~\ref{sec:complete}.}
Two of them differ and both differences are deliberate. First, every method in
this section, incumbents included, is restricted to seeds 0 through 9, because
the Phase 5.5 manifest declared ten seeds rather than twenty. Pooling a
twenty-seed incumbent mean against a ten-seed claimant mean in one row would
compare two designs, so the incumbent columns here are recomputed on the
claimant's seeds and will not always equal the same cell in
Section~\ref{sec:complete}. Second, only D0, D2, D7 and D8 appear in the Stage 1
tables. The screen was declared as a mechanism test on the four regimes that
isolate one mechanism each, so it never ran D1, D3, D4, D5, D6 or D9, and no
Stage 1 table supports a statement about the variants outside those four
regimes. Subsection~\ref{sub:stage2} lifts that restriction for the one variant
that continued, at twenty seeds and on all ten regimes; it remains in force for
\texttt{cwdb\_rmean} and \texttt{cwdb\_xmean}, which were never run elsewhere.
Otherwise the conventions carry over: an entry is a mean over cells with the
standard error in parentheses, a cell is one (grid, regime, $n$, $K$, $M$,
method, seed) combination with arm-specific rows averaged inside the cell first,
bold marks the best method in a regime, and a regime mean pools $n = 500$ and
$n = 1000$, which is why the numbers here are not identical to the per-$n$
figures quoted in the G3.5 screen memo. The Causal-DRF and DRF columns are the
original-code reruns, matching the corrected comparator record of
Section~\ref{sec:complete} rather than the retired local drivers.

\subsection{The contrast surface}

Table~\ref{tab:p55-meanq} is the primary table of this section, because
\texttt{mean\_quantile\_rmse} against \texttt{MEANQ-A-K} is the one metric all
seven methods produce. Four readings are worth stating, and the fourth is the
one the screen memo initially failed to make.

The null and deterministic regimes separate the variants sharply. On D2, where
the true effect is exactly zero, \texttt{rmean} records $0.0191$ against
$0.0551$ for PTA-S and $0.0595$ for R3, and Table~\ref{tab:p55-paired-rmean}
shows those gaps are $8.8$ and $3.2$ paired standard errors. This is the
clearest single result of Phase 5.5: orthogonalisation, not more capacity, is
what removes manufactured heterogeneity, and it removes it against the
comparator that motivated the whole repair track. On D8, the confounded regime
that rule 2 of the original gate treated as primary, \texttt{rmean} beats PTA-S
by $2.8$ standard errors but loses to R3 by $5.9$, so the R-loss buys
null-safety at some cost in confounded accuracy.

D7 is where \texttt{rmean} is weakest, at $0.1033$ against $0.0637$ for PTA-S
and $0.0661$ for R3, losing to every comparator in the table by between $2.2$
and $10.9$ standard errors. D7 is the regime whose effect acts through the shape
of the law rather than its location, so a method that estimates the contrast
directly and never represents the law has the least to work with there. That is
a coherent explanation rather than an excuse, and it is also the reason
\texttt{rmean} is recommended as a benchmark rather than as the claimant.

\texttt{cwdb\_mutau} sits close to R3 everywhere and is not interchangeable with
it. Table~\ref{tab:p55-paired-mutau} shows two regimes where the difference
crosses the decision multiple and both go against \texttt{mutau}: D0 by $3.3$
standard errors and D8 by $2.6$. D2 and D7 are ties at $0.4$ and $0.6$. The
honest summary is that the $\mu/\tau$ leaf rule buys nothing against R3 on the
contrast at these sample sizes while costing a little on two regimes, and that
its case rests on the law and functional tables rather than on this one.

\paragraph{The D0 result, stated rather than passed over.} D0 is the regime with
no treatment effect and no confounding, so it is the closest thing in the suite
to a pure bias check. Every C-WDB variant loses it decisively to PTA-S:
$0.0969$ for \texttt{rmean}, $0.0985$ for \texttt{mutau} and $0.0943$ for R3,
against $0.0544$. In paired terms that is $16.0$ standard errors for
\texttt{rmean} and $37.6$ for \texttt{mutau}. The screen memo originally
reported D0 only as passing the $0.15$ correctness cap while reporting the D2
and D8 differences against the same comparator in full, which is a selective
presentation of a comparison that had already been computed. It is corrected
here and in the memo. The deficit is family-wide, not variant-specific, so it is
best read as a property of the boosted particle representation in the regime
where the truth is a single fixed law and the direct learner has nothing to
estimate. It remains an open weakness of the approach, and Stage 2 should report
D0 alongside D2 and D8 as a comparison rather than as a cap check.

Table~\ref{tab:p55-bary} gives the arm-level barycenter error, which is a
different quantity and is never pooled with the contrast. PTA-S is the best
method in all four regimes there, by a wide margin, which is exactly what a
direct arm-level learner should be good at, and the fact that R3 and
\texttt{mutau} nonetheless win contrasts on D8 is the substantive point: the
contrast is not recoverable from arm-level accuracy alone. Tables
\ref{tab:p55-reftcate} and \ref{tab:p55-refate} report the Wasserstein
reference effects, where \texttt{mutau} is best on D7 for both the conditional
and the marginal version and R3 is best on D2 and D8, all by margins smaller
than the ones separating either from the forest comparators.

\subsection{The law surface}

Only \texttt{cwdb\_mutau} enters Tables \ref{tab:p55-kernel} through
\ref{tab:p55-tail}, and the result there is a near-exact tie with R3.
Table~\ref{tab:p55-paired-mutau-law} makes this precise: on
\texttt{kernel\_law\_error} the paired difference against R3 never exceeds
$1.4$ standard errors in any regime, while the difference against Causal-DRF is
a \texttt{mutau} win in all four at between $8.4$ and $23.7$. This is the
expected result and it is a validity check rather than a discovery. The
$\mu/\tau$ parameterisation provably collapses to the reparameterised shared
tree when the propensity basis is the leaf share, a collapse verified in the
test suite to $10^{-14}$ at both the tree and the model level, so a large law
difference against R3 would have indicated a defect rather than a gain.

The functional tables carry the one substantive law-side finding.
Table~\ref{tab:p55-tcate} shows \texttt{mutau} best or tied-best on
\texttt{grid\_mean} and \texttt{grid\_upper\_tail\_mean} in three of four
regimes, and it shows D7 skewness as a failure region for both C-WDB variants:
$0.1134$ for \texttt{mutau} and $0.0860$ for R3 against $0.0514$ for Causal-DRF
and $0.0580$ for DRF, with the same ordering on the marginal version in
Table~\ref{tab:p55-tate}. This is the same documented failure region
Section~\ref{sec:complete} reports for R3 and it is not repaired by the
$\mu/\tau$ leaf rule; if anything \texttt{mutau} is worse there, which is the
$n = 500$ skewness flag recorded in the screen memo. Pure skewness contrasts
remain a region where the forest comparators are better and no claim should be
made over them.

Table~\ref{tab:p55-mode} is a degenerate check in this section, since D6 is not
in the Phase 5.5 manifest and every regime here is single-mode, so the uniform
$1.0000$ column establishes only that nothing broke.

\subsection{The inherited correctness screen}

Table~\ref{tab:p55-rule1} evaluates inherited rule 1 per sample size, under one
convention applied identically to every method: the D2 false-effect ratio
divides a single $n$'s D2 mean by the best frozen baseline at that same $n$.
Pooling the baseline across sample sizes, or mixing a per-$n$ numerator with a
pooled denominator, produces different numbers and appeared in an earlier draft
of the screen memo; it is not the convention and has been removed there too.

The table carries two entries that need reading with care. The first is
\texttt{cwdb\_xmean}, which fails the D0 cap of $0.15$ outright at both sample
sizes, $0.3545$ and $0.3536$. That is the cleanest available statement of its
failure and it is stronger than the comparative statement the screen memo
originally used. The second is R3 at $n = 500$, which shows a D2 ratio of
$1.26$ against a cap of $1.25$. This is an artefact of the ten-seed restriction
and not a new failure: on the full twenty seeds of its own track the same ratio
is $0.98$ at $n = 500$ and $0.93$ pooled, which is what
Section~\ref{sec:complete} reports and what the gate was evaluated on. It is
printed here rather than suppressed because a reader comparing the two sections
would otherwise find the discrepancy unexplained, and because it is a useful
reminder that a threshold evaluated on half the seeds is a different test.

\subsection{The negative finding, and why it is a mechanism result}

\texttt{cwdb\_xmean} exists to test one hypothesis: that out-of-fold imputation
of the missing potential outcome helps when the arms are unbalanced. That
hypothesis is specific to the imbalance regime, so a result outside it could not
support the method, which is why the \texttt{imbalance} extension grid was
declared before the run. Table~\ref{tab:p55-imbalance} reports it. At a realised
mean propensity of $0.307$, \texttt{xmean} wins D7-imb by $9.9$ paired standard
errors, loses D2-imb by $8.0$, and on D8-imb is behind by $1.9$, which does not
cross the decision multiple. One win in the regime the method was built for,
paid for with a large loss on the null, is not a result that continues to
Stage 2.

The verdict was reached on that table alone. Table~\ref{tab:p55-budget} then
answers a separate question, namely whether the failure is a property of the
mechanism or of the frozen three-step effect budget, and it is included here
because the two answers imply different Stage 2 decisions and the difference
must not be settled by assertion. Two observations settle it. The D0 error is
flat in $n$, $0.3545$ against $0.3536$, which no variance story predicts and an
underfit does; and relaxing the budget from three trees to fifty drops D0 from
$0.3550$ to $0.0658$ while D2 rises from $0.0546$ to $0.1636$ and D8 from
$0.1112$ to $0.2222$ over the very same sweep. So the D0 wall is a bias floor
imposed by a total boosting weight of $0.36$ against the R-loss contrast's
$2.4$, and no budget on the sweep is simultaneously effect-recovering and
null-safe. On the imbalance grid the best available budget is the frozen one and
it still loses D2-imb and D8-imb to \texttt{rmean}. The negative finding is
therefore about the mechanism and survives the objection that the budget was
simply miscalibrated. The structural reason is visible in
Table~\ref{tab:p55-diagnostics}: \texttt{rmean} selects its contrast ridge per
cell by cross-fitted selection, choosing $0$ on D0 and between $50$ and $500$
on the regimes with signal, whereas the X pseudo-outcome regression has only a
global constant and no adaptive regulariser at all.

\subsection{Diagnostics and cost}
\label{sub:p55cost}

Table~\ref{tab:p55-diagnostics} collects the quantities only the Phase 5.5
variants emit. Two entries carry weight. The effective particle support is
$10$ of $10$ in every regime for \texttt{mutau}, against a rule-5 floor of
$0.6M = 6$, so the $\mu/\tau$ leaf rule does not collapse the particle cloud.
The selected ridge is $0$ on D0 for both \texttt{rmean} and \texttt{mutau} and
rises to between $250$ and $430$ on the regimes carrying signal, which is the
adaptive behaviour the negative finding above turns on.

Table~\ref{tab:p55-cost} reports runtime with two ratios rather than one,
because the rule-6 verdict for \texttt{mutau} depends entirely on which
Causal-DRF is the denominator and the project contains two. Against the
original-code rerun on exactly these cells, a median of $14.47$ seconds, the
\texttt{mutau} ratio is $11.8$ and comfortably inside the cap of $60$. Against
the retired local driver at $1.87$ seconds, the ratio is $91.1$ and breaches it.
The $91\times$ figure quoted in the screen memo is the second of these. The
first is the denominator the rest of this report uses, so the defensible
statement is that \texttt{mutau} passes rule 6 against the current comparator
record while remaining the most expensive method in the roster at roughly twice
R3 and twelve times Causal-DRF. Two caveats apply to every row. All Phase 5.5
cells ran on a loaded machine alongside other work, so these are upper bounds
rather than clean measurements; and \texttt{mutau} performs a fully nested
refit for its contrast selection while \texttt{rmean} performs a
nested-approximate one that reuses a single nuisance stack, so the eight-second
median of \texttt{rmean} and the one hundred and seventy second median of
\texttt{mutau} are not measurements of the same procedure.

\subsection{Stage 2: the \texorpdfstring{$\mu/\tau$}{mu/tau} claimant on the whole grid}
\label{sub:stage2}

Everything above is the Stage 1 mechanism screen. Phase 5.5 was declared as a
two-stage design in which the screen decides whether a variant addresses its
intended mechanism and Stage 2 is the frozen comparison against the full
tournament record. Only \texttt{cwdb\_mutau} reached it: \texttt{cwdb\_rmean}
cannot supply a law target by contract and was retained as a benchmark rather
than as a claimant, and \texttt{cwdb\_xmean} failed its screen. Stage 2 runs the
claimant on all ten regimes at the twenty seeds every incumbent already carries,
which is $320$ new cells (the six regimes the screen omitted at twenty seeds,
and a seed top-up on the four it ran at ten). The merge audit returned
\textsc{pass} with $320$ of $320$ cells observed, no duplicate keys and no
failures, against manifest checksum
\texttt{29d7ca09\allowbreak 28aa5420\allowbreak 23a091fd\allowbreak 2ffa3943}.

The stage exists because the screen left the phase's own requirement for this
variant unsatisfied. That requirement names D2, D3, D5, D6, D7, D8 and D9, and
the phase's decision table contains the row stating that a candidate which fixes
D2 by losing D3 or D7 is rejected as the main claimant. D3, D5, D6 and D9 had
never been run for this method, so the missing evidence was exactly the evidence
that could overturn the Stage 1 verdict. Thresholds were frozen in
\texttt{research/simulation\_preregistration\_phase55\_stage2.md} before the
first Stage 2 cell was executed, including a numerical definition of the phase
document's word ``materially'' (a paired \texttt{kernel\_law\_error} loss to R3
crossing two standard errors on D3 or D7) and a declaration of which Causal-DRF
record is the rule 6 denominator, since the ambiguity described in
Subsection~\ref{sub:p55cost} would otherwise have recurred.

\paragraph{Two measurement caveats, both settled before the tables were read.}
The first is that \texttt{mutau}'s implementation was modified after the Stage 1
shards were written, during the adversarial audit of the screen memo, and the
source tree is not under version control, so the change could not be diffed. It
was tested instead, by rerunning three cells at coordinates Stage 1 also ran.
Every accuracy quantity reproduced inside the Stage 1 spread while runtime fell
by a factor near five, so the change is a performance change and the two stages
pool on accuracy but not on cost. A subsequent check across all four shared
regimes and four metrics found no difference between the stages at the five
percent level in sixteen comparisons. Accordingly the accuracy tables pool both
stages and every runtime figure is read from Stage 2 rows alone. The second is
that the pooled paired comparisons in the analysis library keyed their per-cell
values by seed alone, so a comparison pooling both sample sizes silently
overwrote one with the other and used half its design points, with the surviving
half determined by dictionary order rather than by the design. This was found
while reconciling the payload against the generated tables, has been corrected
to key on the full design point, and every number in this subsection is
post-correction. Comparisons restricted to a single sample size were never
affected, which includes the Stage 1 tables above.

\paragraph{The gate.} Table~\ref{tab:p55s2-gate} states each preregistered rule
and what it returned. Rules 1 through 6 all pass. Rule 1 clears both caps with
room, at D0 means of $0.1061$ and $0.0894$ against a cap of $0.15$ and D2
false-effect ratios of $0.88$ and $0.69$ against a cap of $1.25$. Rule 5 passes
with the effective particle support at $10$ of $10$ in every regime against a
floor of six, and with D6 mode coverage at $0.9930$, which is the first
measurement of multimodal behaviour for this variant and is the check the screen
could not perform because D6 was not in its manifest. Rule 6 passes at a ratio
of $13.9$ against a cap of $60$.

\paragraph{The law surface.} Tables~\ref{tab:p55s2-kernel}
through~\ref{tab:p55s2-tail} give the law metrics on all ten regimes and
Table~\ref{tab:p55s2-paired-law} gives the paired differences. Rule 2 passes
comfortably: on \texttt{kernel\_law\_error} \texttt{mutau} beats Causal-DRF in
eight regimes against a requirement of two, ties D9, and loses only D6, where
Causal-DRF records $0.0077$ against \texttt{mutau}'s $0.0244$, a loss of $16.1$
paired standard errors. That D6 deficit is the same multimodality weakness
Section~\ref{sec:complete} documents for the C-WDB family and it is not repaired
by the $\mu/\tau$ leaf rule.

Against C-WDB R3 the law result is a tie in seven regimes of ten and a loss in
three, D3 at $3.1$ standard errors, D5 at $2.4$ and D4 at $2.2$, with no regime
in which \texttt{mutau} is better. The Stage 1 conclusion that the two are
near-indistinguishable on the law therefore survives the widening in most
regimes but not everywhere, and where it breaks it breaks against the claimant.
Against C-WDB V1 the picture is genuinely mixed, with \texttt{mutau} better on
D2, D6, D7 and D8 and worse on D0, D1, D4, D5 and D9, which is the expected
signature of a reparameterisation that helps where the contrast is small
relative to the prognostic surface and costs where it is not.

\paragraph{The contrast surface and the accuracy clause.}
Table~\ref{tab:p55s2-meanq} and Table~\ref{tab:p55s2-paired-meanq} carry the
common contrast target. Against R3, \texttt{mutau} loses D0, D3, D4 and D8 by
between $2.3$ and $3.6$ standard errors, wins D5 by $3.1$, and ties the other
five, so the Stage 1 reading that the $\mu/\tau$ leaf rule buys nothing against
R3 on the contrast holds on the wider grid. Against PTA-S it wins D2, D3 and D8
and loses D0, D1, D4, D5 and D6, with the D0 loss at $55.9$ standard errors,
which is the family-wide D0 weakness recorded above at its most extreme.

The accuracy clause is the one place where Stage 2 improves materially on the
original tournament. That clause, frozen before decisive execution, requires a
full-law candidate to win at least one target that PTA-S also estimates, because
the G3 memo recorded rule 4 passing on three capability wins and \emph{zero}
accuracy wins. \texttt{mutau} returns seventeen accuracy wins spread over seven
regimes (D1, D2, D3, D4, D6, D8 and D9), including the marginal functional on
D6 by $5.9$ standard errors and every reference target on D8. This is a real
gain over the incumbent claimant and it is the strongest argument the variant
has.

\paragraph{The degradation clause, which fails.} Table~\ref{tab:p55s2-degradation}
reports it. On D3, the separate-head-favourable regime that a shared prognostic
and contrast tree puts at risk, \texttt{mutau} loses the law metric to R3 at
$n = 500$ by $3.0$ paired standard errors, and the loss shrinks to $1.3$ at
$n = 1000$. D7 is a tie at both sample sizes, so the pure-shape-transfer half of
the clause is satisfied and the sharing half is not. Under the phase's own
decision table this is the row reading that a candidate which fixes D2 while
losing D3 is rejected as the main claimant and its regularisation tradeoff is
reported instead. The magnitude is small in absolute terms, $0.00078$ against a
level near $0.02$, and it is confined to the smaller sample size; the
preregistration nonetheless defined the clause in standard errors rather than in
absolute terms, before the number existed, and the honest reading is that the
clause fails as written rather than that the threshold should now be
reinterpreted.

\paragraph{Diagnostics and cost.} Table~\ref{tab:p55s2-diagnostics} shows the
effective support at the full ten particles in every regime and the selected
contrast strength moving with the signal, at zero on D0 and D2 and rising on the
regimes that carry an effect, which is the adaptive behaviour the mechanism
depends on. Table~\ref{tab:p55s2-cost} puts the \texttt{mutau} median at
$199.3$ seconds against the original-code Causal-DRF's $14.36$, a ratio of
$13.9$. Every figure there is an upper bound: the Stage 2 cells ran sixteen at a
time on one laptop, and the same implementation on an idle machine finishes a
cell in roughly a fifth of the time. Runtime in this project is dominated by
machine load, so the cost rows order the methods reliably and measure none of
them precisely.

\subsection{What this section supports}

Four statements, and no more than four. First, \texttt{cwdb\_rmean} is a
credible null-safe benchmark: it is the best method in the suite on the null
regime by a wide margin, it beats PTA-S on the confounded regime, and it is
cheap, but it loses D0 and D7 badly enough that it is a benchmark and not a
claimant. Second, \texttt{cwdb\_xmean} is a negative finding on the mechanism,
not a tuning failure, and it produced no Stage 2 manifest.

Third, on the evidence of Stage 2, \texttt{cwdb\_mutau} is a real improvement on
the incumbent claimant in one specific respect and not a replacement for it. The
improvement is the accuracy clause: seventeen wins against PTA-S on targets
PTA-S also estimates, spread over seven regimes, where the original tournament
recorded rule 4 passing on capability wins alone. Alongside it the variant
clears every inherited gate rule, beats Causal-DRF on the law metric in eight
regimes of ten, holds full particle support, and is the first C-WDB variant
measured on D6 multimodality, at a mode coverage of $0.9930$. Against that, it
is never better than R3 on the law metric in any regime and is worse in three,
it loses the common contrast to R3 in four regimes and wins one, and it fails
the preregistered degradation clause on D3 at $n = 500$, which is the row the
phase's decision table maps to rejection as the main claimant. The defensible
position is therefore that the $\mu/\tau$ reparameterisation is a reportable
regularisation tradeoff and a candidate complementary estimator, and that
nothing here licenses replacing R3 with it.

Fourth, the scope limits. The Stage 1 tables speak only to D0, D2, D7 and D8,
and nothing in this section speaks to \texttt{cwdb\_rmean} or
\texttt{cwdb\_xmean} outside those four regimes. Stage 2 covers all ten regimes
for \texttt{cwdb\_mutau} alone. No part of this section speaks to any sample
size other than $500$ and $1000$, to any grid resolution other than $K = 25$, to
any particle count other than $M = 10$, or to uncertainty coverage for any
variant, since none has an interval construction and the estimand contract
forbids substituting a posterior-draw quantity.

\begin{table}[htbp]
\centering
\caption{Grid causal mean error \texttt{mean\_quantile\_rmse} against \texttt{MEANQ-A-K}, main grid, mean over 20 cells ($n\in\{500,1000\}$, ten seeds) with the standard error in parentheses. This is the primary contrast metric and the one every Phase 5.5 variant was screened on.}
\label{tab:p55-meanq}
\scriptsize
\setlength{\tabcolsep}{3.2pt}
\begin{tabular}{@{}lrrrrrrr@{}}
\toprule
Regime & rmean & mutau & xmean & R3 & PTA-S & C-DRF & DRF \\
\midrule
D0 & 0.0969\,(0.0038) & 0.0985\,(0.0022) & 0.3540\,(0.0012) & 0.0943\,(0.0021) & \textbf{0.0544\,(0.0018)} & 0.1637\,(0.0059) & 0.1920\,(0.0061) \\
D2 & \textbf{0.0191\,(0.0030)} & 0.0532\,(0.0098) & 0.0473\,(0.0024) & 0.0595\,(0.0116) & 0.0551\,(0.0039) & 0.0952\,(0.0070) & 0.0743\,(0.0064) \\
D7 & 0.1033\,(0.0040) & 0.0647\,(0.0028) & \textbf{0.0624\,(0.0019)} & 0.0661\,(0.0034) & 0.0637\,(0.0035) & 0.0871\,(0.0066) & 0.0787\,(0.0066) \\
D8 & 0.0885\,(0.0025) & 0.0795\,(0.0031) & 0.0915\,(0.0048) & \textbf{0.0739\,(0.0034)} & 0.0984\,(0.0047) & 0.2513\,(0.0137) & 0.1843\,(0.0084) \\
\bottomrule
\end{tabular}
\end{table}

\begin{table}[htbp]
\centering
\caption{Arm-level barycenter error \texttt{barycenter\_rmse}, main grid. This is an arm-level quantity, never pooled with the contrast of Table~\ref{tab:p55-meanq}.}
\label{tab:p55-bary}
\scriptsize
\setlength{\tabcolsep}{3.2pt}
\begin{tabular}{@{}lrrrrrrr@{}}
\toprule
Regime & rmean & mutau & xmean & R3 & PTA-S & C-DRF & DRF \\
\midrule
D0 & 0.1252\,(0.0028) & 0.1077\,(0.0027) & 0.2605\,(0.0013) & 0.1065\,(0.0026) & \textbf{0.0544\,(0.0019)} & 0.2115\,(0.0055) & 0.2321\,(0.0058) \\
D2 & 0.1490\,(0.0019) & 0.1884\,(0.0053) & 0.1653\,(0.0037) & 0.1908\,(0.0053) & \textbf{0.1294\,(0.0032)} & 0.2509\,(0.0060) & 0.2322\,(0.0047) \\
D7 & 0.1482\,(0.0029) & 0.1892\,(0.0050) & 0.1639\,(0.0041) & 0.1892\,(0.0049) & \textbf{0.1306\,(0.0033)} & 0.2294\,(0.0056) & 0.2375\,(0.0050) \\
D8 & 0.3034\,(0.0030) & 0.2540\,(0.0056) & 0.3285\,(0.0062) & 0.2529\,(0.0061) & \textbf{0.1950\,(0.0054)} & 0.4529\,(0.0105) & 0.4436\,(0.0080) \\
\bottomrule
\end{tabular}
\end{table}

\begin{table}[htbp]
\centering
\caption{Declared primary law metric \texttt{kernel\_law\_error} (squared MMD to the true conditional law), main grid. Only \texttt{mutau} produces a law among the Phase 5.5 variants; \texttt{rmean} and \texttt{xmean} target the contrast only and are absent by construction, not by omission.}
\label{tab:p55-kernel}
\scriptsize
\setlength{\tabcolsep}{3.2pt}
\begin{tabular}{@{}lrrrr@{}}
\toprule
Regime & mutau & R3 & C-DRF & DRF \\
\midrule
D0 & 0.0289\,(0.0016) & \textbf{0.0287\,(0.0017)} & 0.1209\,(0.0050) & 0.1422\,(0.0062) \\
D2 & \textbf{0.0117\,(0.0006)} & 0.0122\,(0.0006) & 0.0188\,(0.0008) & 0.0163\,(0.0006) \\
D7 & 0.0118\,(0.0006) & \textbf{0.0117\,(0.0006)} & 0.0159\,(0.0008) & 0.0170\,(0.0007) \\
D8 & 0.0157\,(0.0007) & \textbf{0.0155\,(0.0007)} & 0.0498\,(0.0021) & 0.0483\,(0.0018) \\
\bottomrule
\end{tabular}
\end{table}

\begin{table}[htbp]
\centering
\caption{Excess energy risk over the true conditional law \texttt{arm\_energy\_risk}, main grid. This is the loss the law methods actually optimise.}
\label{tab:p55-energy}
\scriptsize
\setlength{\tabcolsep}{3.2pt}
\begin{tabular}{@{}lrrrr@{}}
\toprule
Regime & mutau & R3 & C-DRF & DRF \\
\midrule
D0 & 0.0584\,(0.0016) & \textbf{0.0579\,(0.0016)} & 0.1312\,(0.0028) & 0.1458\,(0.0037) \\
D2 & 0.0536\,(0.0026) & 0.0551\,(0.0026) & 0.0456\,(0.0018) & \textbf{0.0407\,(0.0014)} \\
D7 & 0.0545\,(0.0026) & 0.0549\,(0.0026) & \textbf{0.0403\,(0.0016)} & 0.0430\,(0.0015) \\
D8 & 0.0650\,(0.0026) & \textbf{0.0644\,(0.0026)} & 0.1213\,(0.0043) & 0.1187\,(0.0037) \\
\bottomrule
\end{tabular}
\end{table}

\begin{table}[htbp]
\centering
\caption{\texttt{mode\_coverage}, the fraction of true outer modes carrying predicted mass, main grid. Higher is better. D6, the only multimodal regime, is not in the Phase 5.5 manifest, so every row here is a degenerate single-mode check.}
\label{tab:p55-mode}
\scriptsize
\setlength{\tabcolsep}{3.2pt}
\begin{tabular}{@{}lrrrr@{}}
\toprule
Regime & mutau & R3 & C-DRF & DRF \\
\midrule
D0 & \textbf{1.0000\,(0.0000)} & \textbf{1.0000\,(0.0000)} & \textbf{1.0000\,(0.0000)} & \textbf{1.0000\,(0.0000)} \\
D2 & \textbf{1.0000\,(0.0000)} & \textbf{1.0000\,(0.0000)} & \textbf{1.0000\,(0.0000)} & \textbf{1.0000\,(0.0000)} \\
D7 & \textbf{1.0000\,(0.0000)} & \textbf{1.0000\,(0.0000)} & \textbf{1.0000\,(0.0000)} & \textbf{1.0000\,(0.0000)} \\
D8 & \textbf{1.0000\,(0.0000)} & \textbf{1.0000\,(0.0000)} & \textbf{1.0000\,(0.0000)} & \textbf{1.0000\,(0.0000)} \\
\bottomrule
\end{tabular}
\end{table}

\begin{table}[htbp]
\centering
\caption{\texttt{tail\_calibration}, absolute error of a predeclared threshold event under the grid law, main grid.}
\label{tab:p55-tail}
\scriptsize
\setlength{\tabcolsep}{3.2pt}
\begin{tabular}{@{}lrrrr@{}}
\toprule
Regime & mutau & R3 & C-DRF & DRF \\
\midrule
D0 & 0.0701\,(0.0019) & \textbf{0.0691\,(0.0023)} & 0.1794\,(0.0045) & 0.1816\,(0.0058) \\
D2 & 0.1165\,(0.0033) & 0.1165\,(0.0034) & 0.1037\,(0.0024) & \textbf{0.0962\,(0.0024)} \\
D7 & 0.0533\,(0.0017) & 0.0534\,(0.0017) & \textbf{0.0433\,(0.0016)} & 0.0442\,(0.0017) \\
D8 & \textbf{0.0988\,(0.0034)} & 0.1002\,(0.0032) & 0.1451\,(0.0041) & 0.1337\,(0.0039) \\
\bottomrule
\end{tabular}
\end{table}

\begin{table}[htbp]
\centering
\caption{Conditional Wasserstein reference effect \texttt{REF-TCATE-K} with $\nu_\star$ the standard normal, main grid.}
\label{tab:p55-reftcate}
\scriptsize
\setlength{\tabcolsep}{3.2pt}
\begin{tabular}{@{}lrrrrrrr@{}}
\toprule
Regime & rmean & mutau & xmean & R3 & PTA-S & C-DRF & DRF \\
\midrule
D0 & n/a & 0.0550\,(0.0019) & n/a & 0.0496\,(0.0017) & \textbf{0.0174\,(0.0019)} & 0.1033\,(0.0039) & 0.1113\,(0.0077) \\
D2 & n/a & 0.0202\,(0.0035) & n/a & \textbf{0.0184\,(0.0030)} & 0.0235\,(0.0043) & 0.0287\,(0.0045) & 0.0301\,(0.0045) \\
D7 & n/a & \textbf{0.0215\,(0.0027)} & n/a & 0.0232\,(0.0028) & 0.0244\,(0.0035) & 0.0294\,(0.0036) & 0.0296\,(0.0036) \\
D8 & n/a & \textbf{0.0426\,(0.0038)} & n/a & 0.0444\,(0.0036) & 0.0624\,(0.0039) & 0.0695\,(0.0069) & 0.0866\,(0.0080) \\
\bottomrule
\end{tabular}
\end{table}

\begin{table}[htbp]
\centering
\caption{Marginal Wasserstein reference effect \texttt{REF-ATE-K}, main grid.}
\label{tab:p55-refate}
\scriptsize
\setlength{\tabcolsep}{3.2pt}
\begin{tabular}{@{}lrrrrrrr@{}}
\toprule
Regime & rmean & mutau & xmean & R3 & PTA-S & C-DRF & DRF \\
\midrule
D0 & n/a & 0.0436\,(0.0016) & n/a & 0.0385\,(0.0014) & \textbf{0.0047\,(0.0009)} & 0.0129\,(0.0025) & 0.0102\,(0.0031) \\
D2 & n/a & 0.0123\,(0.0036) & n/a & \textbf{0.0096\,(0.0026)} & 0.0169\,(0.0047) & 0.0195\,(0.0051) & 0.0182\,(0.0046) \\
D7 & n/a & \textbf{0.0140\,(0.0029)} & n/a & 0.0150\,(0.0033) & 0.0158\,(0.0040) & 0.0174\,(0.0044) & 0.0167\,(0.0039) \\
D8 & n/a & 0.0239\,(0.0033) & n/a & \textbf{0.0238\,(0.0028)} & 0.0412\,(0.0057) & 0.0491\,(0.0081) & 0.0667\,(0.0101) \\
\bottomrule
\end{tabular}
\end{table}

\begin{table}[htbp]
\centering
\caption{Conditional functional treatment effects \texttt{tcate\_functional\_rmse}, main grid. D7 is the transfer regime the law claim rests on.}
\label{tab:p55-tcate}
\scriptsize
\setlength{\tabcolsep}{3.2pt}
\begin{tabular}{@{}llrrrr@{}}
\toprule
Functional & Regime & mutau & R3 & C-DRF & DRF \\
\midrule
\texttt{grid\_mean} & D0 & 0.0606\,(0.0019) & \textbf{0.0555\,(0.0019)} & 0.1166\,(0.0058) & 0.1093\,(0.0087) \\
 & D2 & \textbf{0.0268\,(0.0042)} & 0.0274\,(0.0043) & 0.0741\,(0.0087) & 0.0527\,(0.0082) \\
 & D7 & \textbf{0.0244\,(0.0032)} & 0.0295\,(0.0040) & 0.0608\,(0.0084) & 0.0535\,(0.0086) \\
 & D8 & \textbf{0.0573\,(0.0049)} & 0.0586\,(0.0037) & 0.2117\,(0.0165) & 0.1306\,(0.0125) \\
\midrule
\texttt{grid\_sd} & D0 & 0.0117\,(0.0009) & \textbf{0.0101\,(0.0008)} & 0.0171\,(0.0014) & 0.0168\,(0.0015) \\
 & D2 & \textbf{0.0169\,(0.0021)} & 0.0173\,(0.0019) & 0.0180\,(0.0022) & 0.0182\,(0.0020) \\
 & D7 & 0.0183\,(0.0016) & 0.0182\,(0.0021) & 0.0198\,(0.0024) & \textbf{0.0182\,(0.0020)} \\
 & D8 & 0.0204\,(0.0022) & \textbf{0.0173\,(0.0022)} & 0.0267\,(0.0021) & 0.0238\,(0.0021) \\
\midrule
\texttt{grid\_skewness} & D0 & 0.0124\,(0.0023) & 0.0120\,(0.0023) & \textbf{0.0000\,(0.0000)} & 0.0000\,(0.0000) \\
 & D2 & 0.0000\,(0.0000) & 0.0000\,(0.0000) & 0.0000\,(0.0000) & \textbf{0.0000\,(0.0000)} \\
 & D7 & 0.1134\,(0.0140) & 0.0860\,(0.0068) & \textbf{0.0514\,(0.0032)} & 0.0580\,(0.0026) \\
 & D8 & 0.0000\,(0.0000) & 0.0000\,(0.0000) & 0.0000\,(0.0000) & \textbf{0.0000\,(0.0000)} \\
\midrule
\texttt{grid\_upper\_tail\_mean} & D0 & 0.0603\,(0.0019) & \textbf{0.0551\,(0.0018)} & 0.1223\,(0.0063) & 0.1129\,(0.0089) \\
 & D2 & \textbf{0.0308\,(0.0040)} & 0.0311\,(0.0045) & 0.0784\,(0.0089) & 0.0567\,(0.0085) \\
 & D7 & \textbf{0.0291\,(0.0033)} & 0.0330\,(0.0038) & 0.0643\,(0.0084) & 0.0572\,(0.0088) \\
 & D8 & 0.0608\,(0.0050) & \textbf{0.0596\,(0.0039)} & 0.2207\,(0.0169) & 0.1361\,(0.0124) \\
\bottomrule
\end{tabular}
\end{table}

\begin{table}[htbp]
\centering
\caption{Marginal functional treatment effects \texttt{tate\_functional\_rmse}, main grid.}
\label{tab:p55-tate}
\scriptsize
\setlength{\tabcolsep}{3.2pt}
\begin{tabular}{@{}llrrrr@{}}
\toprule
Functional & Regime & mutau & R3 & C-DRF & DRF \\
\midrule
\texttt{grid\_mean} & D0 & 0.0064\,(0.0011) & \textbf{0.0052\,(0.0011)} & 0.0575\,(0.0060) & 0.0312\,(0.0048) \\
 & D2 & \textbf{0.0158\,(0.0024)} & 0.0177\,(0.0029) & 0.0657\,(0.0088) & 0.0255\,(0.0056) \\
 & D7 & \textbf{0.0163\,(0.0027)} & 0.0178\,(0.0027) & 0.0462\,(0.0086) & 0.0263\,(0.0057) \\
 & D8 & \textbf{0.0318\,(0.0054)} & 0.0369\,(0.0048) & 0.1850\,(0.0197) & 0.0714\,(0.0136) \\
\midrule
\texttt{grid\_sd} & D0 & 0.0057\,(0.0010) & \textbf{0.0042\,(0.0009)} & 0.0138\,(0.0017) & 0.0132\,(0.0017) \\
 & D2 & \textbf{0.0119\,(0.0021)} & 0.0121\,(0.0017) & 0.0142\,(0.0024) & 0.0137\,(0.0023) \\
 & D7 & \textbf{0.0118\,(0.0014)} & 0.0124\,(0.0019) & 0.0143\,(0.0026) & 0.0136\,(0.0023) \\
 & D8 & 0.0137\,(0.0019) & \textbf{0.0121\,(0.0023)} & 0.0191\,(0.0022) & 0.0192\,(0.0022) \\
\midrule
\texttt{grid\_skewness} & D0 & 0.0102\,(0.0024) & 0.0106\,(0.0024) & \textbf{0.0000\,(0.0000)} & 0.0000\,(0.0000) \\
 & D2 & 0.0000\,(0.0000) & 0.0000\,(0.0000) & 0.0000\,(0.0000) & \textbf{0.0000\,(0.0000)} \\
 & D7 & 0.0977\,(0.0124) & 0.0704\,(0.0059) & 0.0181\,(0.0026) & \textbf{0.0148\,(0.0023)} \\
 & D8 & 0.0000\,(0.0000) & 0.0000\,(0.0000) & \textbf{0.0000\,(0.0000)} & 0.0000\,(0.0000) \\
\midrule
\texttt{grid\_upper\_tail\_mean} & D0 & 0.0071\,(0.0014) & \textbf{0.0055\,(0.0007)} & 0.0677\,(0.0065) & 0.0400\,(0.0052) \\
 & D2 & 0.0204\,(0.0030) & \textbf{0.0197\,(0.0035)} & 0.0706\,(0.0091) & 0.0301\,(0.0066) \\
 & D7 & \textbf{0.0191\,(0.0037)} & 0.0203\,(0.0036) & 0.0504\,(0.0088) & 0.0301\,(0.0065) \\
 & D8 & \textbf{0.0317\,(0.0059)} & 0.0357\,(0.0050) & 0.1953\,(0.0199) & 0.0820\,(0.0141) \\
\bottomrule
\end{tabular}
\end{table}

\begin{table}[htbp]
\centering
\caption{Inherited rule 1, evaluated per sample size on the ten Phase 5.5 seeds. The D0 cap is 0.15, the D2 cap is 0.15, and the D2 false-effect ratio cap is 1.25 times the best frozen baseline at the same $n$.}
\label{tab:p55-rule1}
\small
\begin{tabular}{@{}llrrrrl@{}}
\toprule
Method & $n$ & D0 & D2 & baseline & ratio & Verdict \\
\midrule
rmean & 500 & 0.1095 & 0.0149 & 0.0620 & 0.24 & \textsc{pass} \\
rmean & 1000 & 0.0842 & 0.0233 & 0.0482 & 0.48 & \textsc{pass} \\
\addlinespace
mutau & 500 & 0.1063 & 0.0714 & 0.0620 & 1.15 & \textsc{pass} \\
mutau & 1000 & 0.0907 & 0.0349 & 0.0482 & 0.72 & \textsc{pass} \\
\addlinespace
xmean & 500 & 0.3545 & 0.0553 & 0.0620 & 0.89 & \textbf{FAIL} \\
xmean & 1000 & 0.3536 & 0.0394 & 0.0482 & 0.82 & \textbf{FAIL} \\
\addlinespace
R3 & 500 & 0.1017 & 0.0780 & 0.0620 & 1.26 & \textbf{FAIL} \\
R3 & 1000 & 0.0868 & 0.0410 & 0.0482 & 0.85 & \textsc{pass} \\
\addlinespace
PTA-S & 500 & 0.0623 & 0.0620 & 0.0620 & 1.00 & \textsc{pass} \\
PTA-S & 1000 & 0.0465 & 0.0482 & 0.0482 & 1.00 & \textsc{pass} \\
\addlinespace
C-DRF & 500 & 0.1883 & 0.1086 & 0.0620 & 1.75 & \textbf{FAIL} \\
C-DRF & 1000 & 0.1391 & 0.0818 & 0.0482 & 1.70 & \textbf{FAIL} \\
\addlinespace
DRF & 500 & 0.2168 & 0.0860 & 0.0620 & 1.39 & \textbf{FAIL} \\
DRF & 1000 & 0.1672 & 0.0627 & 0.0482 & 1.30 & \textbf{FAIL} \\
\bottomrule
\end{tabular}
\par\vspace{2pt}\parbox{0.9\textwidth}{\scriptsize The baseline column is the smallest D2 mean among PTA-S, C-WDB-v0, C-WDB-v1, W-DRF-T and Causal-DRF at that $n$, on the same ten seeds. Causal-DRF and DRF are included as rows for reference; they are incumbents, not claimants, and rule 1 was never a gate on them.}
\end{table}

\begin{table}[htbp]
\centering
\caption{Seed-paired \texttt{mean\_quantile\_rmse} differences for \texttt{cwdb\_rmean}. The D8 column is the primary confounding test and the D0 column is the one the screen memo initially omitted.}
\label{tab:p55-paired-rmean}
\scriptsize
\setlength{\tabcolsep}{4pt}
\begin{tabular}{@{}lrrrrrrrrr@{}}
\toprule
Regime & \multicolumn{3}{c}{vs PTA-S} & \multicolumn{3}{c}{vs R3} & \multicolumn{3}{c}{vs C-DRF} \\
  & diff. & SE & $|$diff$|/$SE & diff. & SE & $|$diff$|/$SE & diff. & SE & $|$diff$|/$SE \\
\midrule
D0 & +0.0425$^{\dagger}$ & 0.0026 & 16.0 & +0.0026 & 0.0033 & 0.8 & -0.0669$^{\ast}$ & 0.0039 & 17.0 \\
D2 & -0.0360$^{\ast}$ & 0.0041 & 8.8 & -0.0404$^{\ast}$ & 0.0127 & 3.2 & -0.0761$^{\ast}$ & 0.0082 & 9.3 \\
D7 & +0.0396$^{\dagger}$ & 0.0036 & 10.9 & +0.0372$^{\dagger}$ & 0.0039 & 9.5 & +0.0163$^{\dagger}$ & 0.0073 & 2.2 \\
D8 & -0.0100$^{\ast}$ & 0.0036 & 2.8 & +0.0145$^{\dagger}$ & 0.0025 & 5.9 & -0.1628$^{\ast}$ & 0.0134 & 12.1 \\
\bottomrule
\end{tabular}
\par\vspace{2pt}\parbox{0.96\textwidth}{\scriptsize Negative favours the claimant. $\ast$ marks a claimant win and $\dagger$ a comparator win, both at more than two paired standard errors, which is the frozen decision multiple. The third column of each block is the same quantity expressed in standard errors, so a reader can see how far past the threshold a result sits rather than only that it crossed.}
\end{table}

\begin{table}[htbp]
\centering
\caption{Seed-paired \texttt{mean\_quantile\_rmse} differences for \texttt{cwdb\_mutau}. The R3 block is the informative one: \texttt{mutau} is a reparameterisation of the same shared tree, so a difference against R3 is a difference the $\mu/\tau$ leaf rule caused.}
\label{tab:p55-paired-mutau}
\scriptsize
\setlength{\tabcolsep}{4pt}
\begin{tabular}{@{}lrrrrrrrrr@{}}
\toprule
Regime & \multicolumn{3}{c}{vs R3} & \multicolumn{3}{c}{vs PTA-S} & \multicolumn{3}{c}{vs C-DRF} \\
  & diff. & SE & $|$diff$|/$SE & diff. & SE & $|$diff$|/$SE & diff. & SE & $|$diff$|/$SE \\
\midrule
D0 & +0.0042$^{\dagger}$ & 0.0013 & 3.3 & +0.0441$^{\dagger}$ & 0.0012 & 37.6 & -0.0652$^{\ast}$ & 0.0045 & 14.5 \\
D2 & -0.0063 & 0.0148 & 0.4 & -0.0019 & 0.0084 & 0.2 & -0.0420$^{\ast}$ & 0.0120 & 3.5 \\
D7 & -0.0015 & 0.0027 & 0.6 & +0.0009 & 0.0024 & 0.4 & -0.0224$^{\ast}$ & 0.0060 & 3.7 \\
D8 & +0.0055$^{\dagger}$ & 0.0022 & 2.6 & -0.0190$^{\ast}$ & 0.0042 & 4.6 & -0.1718$^{\ast}$ & 0.0143 & 12.1 \\
\bottomrule
\end{tabular}
\par\vspace{2pt}\parbox{0.96\textwidth}{\scriptsize Negative favours the claimant. $\ast$ marks a claimant win and $\dagger$ a comparator win, both at more than two paired standard errors, which is the frozen decision multiple. The third column of each block is the same quantity expressed in standard errors, so a reader can see how far past the threshold a result sits rather than only that it crossed.}
\end{table}

\begin{table}[htbp]
\centering
\caption{Seed-paired \texttt{mean\_quantile\_rmse} differences for \texttt{cwdb\_xmean} on the main grid. The imbalance grid, which is where its premise lives, is Table~\ref{tab:p55-imbalance}.}
\label{tab:p55-paired-xmean}
\scriptsize
\setlength{\tabcolsep}{4pt}
\begin{tabular}{@{}lrrrrrrrrr@{}}
\toprule
Regime & \multicolumn{3}{c}{vs rmean} & \multicolumn{3}{c}{vs PTA-S} & \multicolumn{3}{c}{vs R3} \\
  & diff. & SE & $|$diff$|/$SE & diff. & SE & $|$diff$|/$SE & diff. & SE & $|$diff$|/$SE \\
\midrule
D0 & +0.2572$^{\dagger}$ & 0.0037 & 70.3 & +0.2996$^{\dagger}$ & 0.0021 & 140.3 & +0.2597$^{\dagger}$ & 0.0027 & 97.7 \\
D2 & +0.0283$^{\dagger}$ & 0.0042 & 6.7 & -0.0077$^{\ast}$ & 0.0031 & 2.5 & -0.0122 & 0.0118 & 1.0 \\
D7 & -0.0410$^{\ast}$ & 0.0037 & 11.1 & -0.0014 & 0.0029 & 0.5 & -0.0038 & 0.0025 & 1.5 \\
D8 & +0.0031 & 0.0049 & 0.6 & -0.0069 & 0.0056 & 1.2 & +0.0176$^{\dagger}$ & 0.0036 & 4.9 \\
\bottomrule
\end{tabular}
\par\vspace{2pt}\parbox{0.96\textwidth}{\scriptsize Negative favours the claimant. $\ast$ marks a claimant win and $\dagger$ a comparator win, both at more than two paired standard errors, which is the frozen decision multiple. The third column of each block is the same quantity expressed in standard errors, so a reader can see how far past the threshold a result sits rather than only that it crossed.}
\end{table}

\begin{table}[htbp]
\centering
\caption{Seed-paired \texttt{kernel\_law\_error} differences for \texttt{cwdb\_mutau}, the law metric.}
\label{tab:p55-paired-mutau-law}
\scriptsize
\setlength{\tabcolsep}{4pt}
\begin{tabular}{@{}lrrrrrr@{}}
\toprule
Regime & \multicolumn{3}{c}{vs R3} & \multicolumn{3}{c}{vs C-DRF} \\
  & diff. & SE & $|$diff$|/$SE & diff. & SE & $|$diff$|/$SE \\
\midrule
D0 & +0.0001 & 0.0006 & 0.2 & -0.0920$^{\ast}$ & 0.0039 & 23.7 \\
D2 & -0.0005 & 0.0004 & 1.4 & -0.0071$^{\ast}$ & 0.0006 & 11.6 \\
D7 & +0.0000 & 0.0002 & 0.1 & -0.0041$^{\ast}$ & 0.0005 & 8.4 \\
D8 & +0.0003 & 0.0003 & 1.0 & -0.0341$^{\ast}$ & 0.0016 & 21.3 \\
\bottomrule
\end{tabular}
\par\vspace{2pt}\parbox{0.96\textwidth}{\scriptsize Negative favours the claimant. $\ast$ marks a claimant win and $\dagger$ a comparator win, both at more than two paired standard errors, which is the frozen decision multiple. The third column of each block is the same quantity expressed in standard errors, so a reader can see how far past the threshold a result sits rather than only that it crossed.}
\end{table}

\begin{table}[htbp]
\centering
\caption{The \texttt{imbalance} extension grid, $n=500$, ten seeds. This grid was added for WP5.5-D: the vector X-learner's claimed advantage is specific to arm imbalance, so a result outside this grid could not support it. The last two columns are the seed-paired difference of \texttt{xmean} against \texttt{rmean} on \texttt{mean\_quantile\_rmse}.}
\label{tab:p55-imbalance}
\small
\begin{tabular}{@{}lrrrrrr@{}}
\toprule
& \multicolumn{2}{c}{\texttt{mean\_quantile\_rmse}} & \multicolumn{2}{c}{\texttt{barycenter\_rmse}} & \multicolumn{2}{c}{paired diff.} \\
\cmidrule(lr){2-3}\cmidrule(lr){4-5}\cmidrule(lr){6-7}
Regime & rmean & xmean & rmean & xmean & diff. & SE \\
\midrule
D2-imb & 0.0153\,(0.0056) & 0.0706\,(0.0053) & 0.1585\,(0.0028) & 0.1797\,(0.0026) & +0.0553$^{\dagger}$ & 0.0069 \\
D7-imb & 0.1289\,(0.0039) & 0.0833\,(0.0043) & 0.1722\,(0.0030) & 0.1785\,(0.0039) & -0.0456$^{\ast}$ & 0.0046 \\
D8-imb & 0.0911\,(0.0026) & 0.1033\,(0.0054) & 0.2954\,(0.0042) & 0.3103\,(0.0055) & +0.0121 & 0.0064 \\
\bottomrule
\end{tabular}
\par\vspace{2pt}\parbox{0.9\textwidth}{\scriptsize Negative favours \texttt{xmean}; $\dagger$ marks an \texttt{rmean} win beyond two paired standard errors. The imbalance actually realised, as the mean fitted propensity: D2-imb 0.307, D7-imb 0.307, D8-imb 0.307.}
\end{table}

\begin{table}[htbp]
\centering
\caption{Post-decision diagnostic for \texttt{cwdb\_xmean}: \texttt{mean\_quantile\_rmse} as the frozen three-step effect budget is relaxed, five seeds. The frozen budget is the first column. No column is simultaneously effect-recovering and null-safe, and on the imbalance grid the best available column is still the frozen one and still loses to \texttt{rmean}.}
\label{tab:p55-budget}
\small
\begin{tabular}{@{}lrrrrr@{}}
\toprule
Cell & $B=3$ & $B=10$ & $B=20$ & $B=50$ & \texttt{rmean} \\
\midrule
\texttt{D0/n500} & 0.3550 & 0.1669 & 0.0850 & 0.0658 & n/a \\
\texttt{D0/n1000} & 0.3510 & 0.1644 & 0.0767 & 0.0520 & n/a \\
\texttt{D2/n500} & 0.0546 & 0.0925 & 0.1217 & 0.1636 & n/a \\
\texttt{D2/n1000} & 0.0387 & 0.0716 & 0.0958 & 0.1295 & n/a \\
\texttt{D8/n500} & 0.1112 & 0.1472 & 0.1770 & 0.2222 & n/a \\
\texttt{D8/n1000} & 0.0854 & 0.1145 & 0.1446 & 0.1810 & n/a \\
\midrule
\texttt{D2-imb/n500} & 0.0682 & 0.1132 & 0.1499 & 0.1926 & 0.0086 \\
\texttt{D7-imb/n500} & 0.0816 & 0.1181 & 0.1483 & 0.1893 & 0.1305 \\
\texttt{D8-imb/n500} & 0.1017 & 0.1532 & 0.1900 & 0.2338 & 0.0931 \\
\bottomrule
\end{tabular}
\par\vspace{2pt}\parbox{0.9\textwidth}{\scriptsize $B$ is \texttt{n\_estimators} in the effect regression; the learning rate stays at 0.12, so the frozen budget carries total boosting weight 0.36 against the R-loss contrast's 2.4. Generated by \texttt{research/checks/phase55\_xmean\_budget\_probe.py}. This probe explains a verdict already reached and no threshold depends on it.}
\end{table}

\begin{table}[htbp]
\centering
\caption{Mechanism diagnostics recorded by the Phase 5.5 variants, main grid, mean over the twenty cells of a regime with the standard error in parentheses. These are not comparisons: no incumbent emits them.}
\label{tab:p55-diagnostics}
\scriptsize
\setlength{\tabcolsep}{4pt}
\begin{tabular}{@{}llrrrr@{}}
\toprule
Method & Diagnostic & D0 & D2 & D7 & D8 \\
\midrule
rmean & selected ridge $\lambda$ & 0.000\,(0.000) & 432.500\,(36.863) & 50.000\,(0.000) & 297.500\,(51.360) \\
mutau & selected ridge $\lambda$ & 0.000\,(0.000) & 272.500\,(52.249) & 252.500\,(51.360) & 297.500\,(51.360) \\
mutau & boosting steps used & 100.000\,(0.000) & 100.000\,(0.000) & 100.000\,(0.000) & 100.000\,(0.000) \\
mutau & effective particle support & 10.000\,(0.000) & 10.000\,(0.000) & 10.000\,(0.000) & 10.000\,(0.000) \\
xmean & $\bar{\hat e}$ & 0.497\,(0.005) & 0.497\,(0.005) & 0.497\,(0.005) & 0.496\,(0.004) \\
xmean & sd of $\hat e$ & 0.128\,(0.004) & 0.128\,(0.004) & 0.128\,(0.004) & 0.295\,(0.003) \\
rmean & training risk & 0.011\,(0.000) & 0.189\,(0.002) & 0.185\,(0.002) & 0.255\,(0.004) \\
mutau & training risk & 0.023\,(0.000) & 0.130\,(0.006) & 0.131\,(0.006) & 0.158\,(0.005) \\
\bottomrule
\end{tabular}
\par\vspace{2pt}\parbox{0.94\textwidth}{\scriptsize The effective support is the number of the $M=10$ particles carrying non-negligible weight, so 10 is the ceiling and the floor for rule 5 is $0.6M=6$. The selected $\lambda$ is chosen per cell from $\{0, 50, 500\}$ by cross-fitted selection, so its mean is a per-regime average of a discrete choice and not itself a tuned value.}
\end{table}

\begin{table}[htbp]
\centering
\caption{Operational cost on the Phase 5.5 coordinates: every \texttt{main} grid cell of D0, D2, D7 and D8 at ten seeds. Two cost ratios are reported against the inherited cap of 60, because the rule-6 verdict depends on which Causal-DRF is the denominator and the project contains two.}
\label{tab:p55-cost}
\small
\begin{tabular}{@{}lrrrrrrr@{}}
\toprule
Method & Cells & Median (s) & Mean (s) & Max (s) & Ratio$_{\text{orig}}$ & Ratio$_{\text{local}}$ & Median RAM (MB) \\
\midrule
C-WDB rmean (vector R-learner) & 80 & 8.27 & 8.77 & 15.8 & 0.6 & 4.4 & 0 \\
C-WDB mutau ($\mu/\tau$ shared tree) & 80 & 170.19 & 180.68 & 321.4 & 11.8 & 91.1 & 0 \\
C-WDB xmean (vector X-learner) & 80 & 5.03 & 5.40 & 9.8 & 0.3 & 2.7 & 0 \\
C-WDB R3 (cv-ridge) & 80 & 84.47 & 87.11 & 134.0 & 5.8 & 45.2 & 0 \\
PTA-S & 80 & 35.36 & 35.73 & 51.4 & 2.4 & 18.9 & 0 \\
Causal-DRF & 80 & 14.47 & 13.10 & 38.1 & 1.0 & 7.7 & 336 \\
DRF (separate, paper benchmark) & 80 & 11.32 & 10.93 & 33.4 & 0.8 & 6.1 & 431 \\
\bottomrule
\end{tabular}
\par\vspace{2pt}\parbox{0.94\textwidth}{\scriptsize Ratio$_{\text{orig}}$ divides by the original-code Causal-DRF rerun on exactly these cells (median 14.47\,s), which is the comparator record the rest of this report uses. Ratio$_{\text{local}}$ divides by the retired local Causal-DRF driver pooled over its own \texttt{main} grid, all ten regimes (median 1.87\,s), which is the denominator the G3.5 screen memo quoted. Every Phase 5.5 cell ran on a loaded machine alongside other work, so these are upper bounds on the variants' cost rather than clean measurements; the ratios are reported anyway because the cap is a declared gate and a favourable rerun cannot be assumed.}
\end{table}

\begin{table}[htbp]
\centering
\caption{Phase 5.5 Stage 2 against the rules frozen in \texttt{research/simulation\_preregistration\_phase55\_stage2.md} before the first Stage 2 cell was executed. Claimant \texttt{cwdb\_mutau}, ten regimes, twenty seeds.}
\label{tab:p55s2-gate}
\scriptsize
\setlength{\tabcolsep}{4pt}
\begin{tabular}{@{}p{0.30\textwidth}p{0.50\textwidth}l@{}}
\toprule
Rule & Result & Verdict \\
\midrule
Rule 1, correctness and nulls & D0 0.1061 and 0.0894 under 0.15; D2 ratio 0.88 and 0.69 under 1.25 & \textbf{PASS} \\
Rule 2, law advantage over Causal-DRF & 8 regimes of 2 required: D0, D1, D2, D3, D4, D5, D7, D8 & \textbf{PASS} \\
Rules 3 and 4, transfer and direct learner & 25 functional or reference targets won against Causal-DRF; 17 accuracy wins against PTA-S across 7 regimes & \textbf{PASS} \\
Rule 5, no particle collapse & effective support 10.0 of 10 against a floor of 6; D6 mode coverage 0.9930 & \textbf{PASS} \\
Rule 6, cost & median 199.3\,s against 14.36\,s, ratio 13.9 against a cap of 60 & \textbf{PASS} \\
Accuracy clause, at least one win over PTA-S & 17 wins on targets both methods estimate, in D1, D2, D3, D4, D6, D8, D9 & \textbf{PASS} \\
Degradation clause, D3 sharing and D7 shape & D3@n500 crosses the multiple against R3 & \textbf{FAIL} \\
\bottomrule
\end{tabular}
\par\vspace{2pt}\parbox{0.96\textwidth}{\scriptsize Rules 1 through 6 are the inherited G3 gate rules. The last two rows are the clauses the phase document froze specifically for a Stage 2 full-law candidate. The degradation clause was given its numerical definition (a paired \texttt{kernel\_law\_error} loss to R3 crossing two standard errors on D3 or D7) in the preregistration, before the result existed.}
\end{table}

\begin{table}[htbp]
\centering
\caption{Stage 2 grid causal mean: \texttt{mean\_quantile\_rmse} against \texttt{MEANQ-A-K}, every regime, twenty seeds. Mean over cells with the standard error in parentheses; bold marks the best method in a regime.}
\label{tab:p55s2-meanq}
\scriptsize
\setlength{\tabcolsep}{3.2pt}
\begin{tabular}{@{}lrrrrrr@{}}
\toprule
Regime & mutau & R3 & v1 & PTA-S & C-DRF & DRF \\
\midrule
D0 & 0.0977\,(0.0015) & 0.0952\,(0.0016) & 0.1091\,(0.0021) & \textbf{0.0547\,(0.0013)} & 0.1621\,(0.0038) & 0.1935\,(0.0048) \\
D1 & 0.2053\,(0.0044) & 0.2043\,(0.0042) & 0.1899\,(0.0042) & \textbf{0.1427\,(0.0030)} & 0.2119\,(0.0062) & 0.2017\,(0.0049) \\
D2 & \textbf{0.0433\,(0.0054)} & 0.0506\,(0.0063) & 0.1466\,(0.0038) & 0.0545\,(0.0030) & 0.0960\,(0.0046) & 0.0754\,(0.0042) \\
D3 & 0.2896\,(0.0054) & \textbf{0.2861\,(0.0049)} & 0.2891\,(0.0058) & 0.2947\,(0.0058) & 0.4301\,(0.0060) & 0.4088\,(0.0072) \\
D4 & 0.2062\,(0.0053) & 0.2002\,(0.0048) & 0.1979\,(0.0051) & \textbf{0.1372\,(0.0036)} & 0.1603\,(0.0062) & 0.1697\,(0.0078) \\
D5 & 0.1346\,(0.0131) & 0.1644\,(0.0133) & 0.1757\,(0.0037) & \textbf{0.0500\,(0.0031)} & 0.0885\,(0.0038) & 0.1004\,(0.0041) \\
D6 & 0.2530\,(0.0120) & 0.2477\,(0.0108) & 0.2898\,(0.0125) & 0.2320\,(0.0098) & 0.2567\,(0.0128) & \textbf{0.2247\,(0.0129)} \\
D7 & 0.0713\,(0.0056) & 0.0683\,(0.0042) & 0.1460\,(0.0037) & \textbf{0.0636\,(0.0026)} & 0.0870\,(0.0042) & 0.0794\,(0.0043) \\
D8 & 0.0776\,(0.0020) & \textbf{0.0735\,(0.0022)} & 0.1848\,(0.0041) & 0.0950\,(0.0030) & 0.2638\,(0.0104) & 0.1790\,(0.0053) \\
D9 & 0.2819\,(0.0094) & 0.2819\,(0.0087) & 0.2503\,(0.0065) & \textbf{0.1900\,(0.0045)} & 0.2891\,(0.0095) & 0.3034\,(0.0101) \\
\bottomrule
\end{tabular}
\end{table}

\begin{table}[htbp]
\centering
\caption{Stage 2 arm-level barycenter error. A different quantity from the contrast and never pooled with it.}
\label{tab:p55s2-bary}
\scriptsize
\setlength{\tabcolsep}{3.2pt}
\begin{tabular}{@{}lrrrrrr@{}}
\toprule
Regime & mutau & R3 & v1 & PTA-S & C-DRF & DRF \\
\midrule
D0 & 0.1077\,(0.0018) & 0.1062\,(0.0019) & 0.1027\,(0.0016) & \textbf{0.0544\,(0.0014)} & 0.2098\,(0.0035) & 0.2313\,(0.0040) \\
D1 & 0.2177\,(0.0040) & 0.2174\,(0.0041) & 0.2124\,(0.0038) & \textbf{0.1497\,(0.0028)} & 0.2632\,(0.0064) & 0.2448\,(0.0043) \\
D2 & 0.1871\,(0.0035) & 0.1886\,(0.0035) & 0.1974\,(0.0037) & \textbf{0.1316\,(0.0027)} & 0.2506\,(0.0041) & 0.2302\,(0.0033) \\
D3 & 0.2362\,(0.0039) & 0.2340\,(0.0037) & 0.2357\,(0.0041) & \textbf{0.1950\,(0.0033)} & 0.3028\,(0.0042) & 0.2881\,(0.0051) \\
D4 & 0.2234\,(0.0042) & 0.2206\,(0.0040) & 0.2170\,(0.0040) & \textbf{0.1402\,(0.0030)} & 0.2236\,(0.0065) & 0.2289\,(0.0085) \\
D5 & 0.1808\,(0.0050) & 0.1846\,(0.0049) & 0.1846\,(0.0034) & \textbf{0.1405\,(0.0030)} & 0.2167\,(0.0037) & 0.2346\,(0.0034) \\
D6 & 0.3752\,(0.0087) & 0.3731\,(0.0086) & 0.3789\,(0.0096) & 0.2788\,(0.0051) & \textbf{0.2764\,(0.0070)} & 0.2977\,(0.0090) \\
D7 & 0.1889\,(0.0037) & 0.1884\,(0.0035) & 0.1971\,(0.0038) & \textbf{0.1327\,(0.0027)} & 0.2285\,(0.0039) & 0.2352\,(0.0035) \\
D8 & 0.2498\,(0.0040) & 0.2500\,(0.0042) & 0.2712\,(0.0037) & \textbf{0.1943\,(0.0034)} & 0.4517\,(0.0075) & 0.4384\,(0.0051) \\
D9 & 0.2589\,(0.0064) & 0.2596\,(0.0059) & 0.2417\,(0.0051) & \textbf{0.1773\,(0.0037)} & 0.2692\,(0.0058) & 0.3070\,(0.0079) \\
\bottomrule
\end{tabular}
\end{table}

\begin{table}[htbp]
\centering
\caption{Stage 2 Wasserstein reference effect, \texttt{REF-ATE-K}.}
\label{tab:p55s2-refate}
\scriptsize
\setlength{\tabcolsep}{3.2pt}
\begin{tabular}{@{}lrrrrrr@{}}
\toprule
Regime & mutau & R3 & v1 & PTA-S & C-DRF & DRF \\
\midrule
D0 & 0.0428\,(0.0011) & 0.0393\,(0.0010) & 0.0421\,(0.0017) & \textbf{0.0058\,(0.0008)} & 0.0126\,(0.0015) & 0.0114\,(0.0018) \\
D1 & 0.0420\,(0.0052) & 0.0381\,(0.0049) & 0.0269\,(0.0035) & \textbf{0.0188\,(0.0029)} & 0.0206\,(0.0030) & 0.0211\,(0.0028) \\
D2 & 0.0100\,(0.0019) & \textbf{0.0099\,(0.0016)} & 0.0205\,(0.0027) & 0.0168\,(0.0026) & 0.0178\,(0.0029) & 0.0175\,(0.0026) \\
D3 & 0.0213\,(0.0028) & \textbf{0.0209\,(0.0028)} & 0.0301\,(0.0034) & 0.0251\,(0.0030) & 0.0238\,(0.0032) & 0.0241\,(0.0031) \\
D4 & 0.0710\,(0.0062) & 0.0562\,(0.0049) & 0.0298\,(0.0039) & \textbf{0.0198\,(0.0029)} & 0.0240\,(0.0033) & 0.0251\,(0.0029) \\
D5 & 0.1210\,(0.0039) & 0.1152\,(0.0036) & 0.0998\,(0.0043) & 0.0223\,(0.0024) & 0.0229\,(0.0026) & \textbf{0.0214\,(0.0026)} \\
D6 & \textbf{0.0240\,(0.0031)} & 0.0262\,(0.0030) & 0.0643\,(0.0072) & 0.0424\,(0.0064) & 0.0422\,(0.0061) & 0.0417\,(0.0067) \\
D7 & \textbf{0.0134\,(0.0017)} & 0.0143\,(0.0019) & 0.0204\,(0.0023) & 0.0158\,(0.0022) & 0.0164\,(0.0025) & 0.0167\,(0.0022) \\
D8 & 0.0250\,(0.0024) & \textbf{0.0233\,(0.0021)} & 0.0328\,(0.0040) & 0.0397\,(0.0042) & 0.0458\,(0.0057) & 0.0645\,(0.0066) \\
D9 & 0.1296\,(0.0103) & 0.1340\,(0.0102) & 0.0437\,(0.0056) & 0.0449\,(0.0049) & \textbf{0.0250\,(0.0028)} & 0.0267\,(0.0034) \\
\bottomrule
\end{tabular}
\end{table}

\begin{table}[htbp]
\centering
\caption{Stage 2 conditional Wasserstein reference effect, \texttt{REF-TCATE-K}.}
\label{tab:p55s2-reftcate}
\scriptsize
\setlength{\tabcolsep}{3.2pt}
\begin{tabular}{@{}lrrrrrr@{}}
\toprule
Regime & mutau & R3 & v1 & PTA-S & C-DRF & DRF \\
\midrule
D0 & 0.0542\,(0.0013) & 0.0508\,(0.0012) & 0.0575\,(0.0016) & \textbf{0.0158\,(0.0012)} & 0.1037\,(0.0026) & 0.1140\,(0.0055) \\
D1 & 0.0775\,(0.0056) & 0.0734\,(0.0054) & \textbf{0.0562\,(0.0032)} & 0.0578\,(0.0039) & 0.1047\,(0.0027) & 0.1138\,(0.0039) \\
D2 & \textbf{0.0168\,(0.0019)} & 0.0170\,(0.0018) & 0.0328\,(0.0024) & 0.0232\,(0.0024) & 0.0270\,(0.0027) & 0.0291\,(0.0026) \\
D3 & 0.0384\,(0.0027) & 0.0375\,(0.0030) & 0.0441\,(0.0034) & 0.0354\,(0.0024) & 0.0355\,(0.0028) & \textbf{0.0337\,(0.0026)} \\
D4 & 0.1132\,(0.0073) & 0.0936\,(0.0057) & 0.0692\,(0.0039) & \textbf{0.0607\,(0.0046)} & 0.1098\,(0.0044) & 0.1193\,(0.0074) \\
D5 & 0.1234\,(0.0038) & 0.1180\,(0.0035) & 0.1032\,(0.0042) & \textbf{0.0286\,(0.0025)} & 0.0330\,(0.0027) & 0.0311\,(0.0024) \\
D6 & \textbf{0.0355\,(0.0031)} & 0.0389\,(0.0032) & 0.0897\,(0.0070) & 0.0549\,(0.0062) & 0.0646\,(0.0058) & 0.0610\,(0.0062) \\
D7 & \textbf{0.0208\,(0.0016)} & 0.0215\,(0.0017) & 0.0309\,(0.0021) & 0.0243\,(0.0019) & 0.0278\,(0.0022) & 0.0285\,(0.0020) \\
D8 & 0.0421\,(0.0025) & \textbf{0.0406\,(0.0026)} & 0.0576\,(0.0038) & 0.0558\,(0.0038) & 0.0642\,(0.0050) & 0.0830\,(0.0056) \\
D9 & 0.1675\,(0.0096) & 0.1708\,(0.0096) & \textbf{0.0832\,(0.0058)} & 0.1021\,(0.0048) & 0.1044\,(0.0031) & 0.1277\,(0.0036) \\
\bottomrule
\end{tabular}
\end{table}

\begin{table}[htbp]
\centering
\caption{Stage 2 primary law metric, \texttt{kernel\_law\_error}. Only methods that produce a conditional law appear.}
\label{tab:p55s2-kernel}
\scriptsize
\setlength{\tabcolsep}{3.2pt}
\begin{tabular}{@{}lrrrrrr@{}}
\toprule
Regime & mutau & R3 & v1 & W-DRF-T & C-DRF & DRF \\
\midrule
D0 & 0.0284\,(0.0011) & 0.0280\,(0.0012) & \textbf{0.0252\,(0.0010)} & 0.0964\,(0.0030) & 0.1166\,(0.0035) & 0.1365\,(0.0044) \\
D1 & 0.0138\,(0.0005) & 0.0137\,(0.0005) & 0.0133\,(0.0005) & \textbf{0.0123\,(0.0003)} & 0.0191\,(0.0009) & 0.0171\,(0.0006) \\
D2 & \textbf{0.0114\,(0.0004)} & 0.0115\,(0.0004) & 0.0127\,(0.0005) & 0.0125\,(0.0003) & 0.0184\,(0.0006) & 0.0158\,(0.0005) \\
D3 & 0.0168\,(0.0006) & 0.0162\,(0.0005) & 0.0165\,(0.0006) & \textbf{0.0158\,(0.0006)} & 0.0259\,(0.0007) & 0.0238\,(0.0008) \\
D4 & 0.0118\,(0.0004) & 0.0115\,(0.0004) & 0.0111\,(0.0004) & \textbf{0.0065\,(0.0004)} & 0.0128\,(0.0007) & 0.0141\,(0.0010) \\
D5 & 0.0250\,(0.0009) & 0.0234\,(0.0006) & \textbf{0.0234\,(0.0007)} & 0.0363\,(0.0010) & 0.0431\,(0.0013) & 0.0478\,(0.0012) \\
D6 & 0.0244\,(0.0012) & 0.0246\,(0.0013) & 0.0278\,(0.0016) & \textbf{0.0066\,(0.0003)} & 0.0077\,(0.0004) & 0.0082\,(0.0005) \\
D7 & 0.0115\,(0.0004) & \textbf{0.0114\,(0.0004)} & 0.0127\,(0.0005) & 0.0129\,(0.0003) & 0.0155\,(0.0005) & 0.0165\,(0.0005) \\
D8 & \textbf{0.0151\,(0.0005)} & 0.0151\,(0.0005) & 0.0176\,(0.0005) & 0.0376\,(0.0008) & 0.0489\,(0.0014) & 0.0466\,(0.0011) \\
D9 & 0.0199\,(0.0010) & 0.0197\,(0.0009) & 0.0169\,(0.0007) & \textbf{0.0162\,(0.0006)} & 0.0200\,(0.0008) & 0.0260\,(0.0013) \\
\bottomrule
\end{tabular}
\end{table}

\begin{table}[htbp]
\centering
\caption{Stage 2 excess energy risk by arm, averaged inside the cell.}
\label{tab:p55s2-energy}
\scriptsize
\setlength{\tabcolsep}{3.2pt}
\begin{tabular}{@{}lrrrrrr@{}}
\toprule
Regime & mutau & R3 & v1 & W-DRF-T & C-DRF & DRF \\
\midrule
D0 & 0.0594\,(0.0012) & 0.0586\,(0.0012) & \textbf{0.0562\,(0.0010)} & 0.1192\,(0.0017) & 0.1318\,(0.0019) & 0.1464\,(0.0026) \\
D1 & 0.0626\,(0.0020) & 0.0624\,(0.0021) & 0.0624\,(0.0020) & \textbf{0.0360\,(0.0008)} & 0.0507\,(0.0019) & 0.0477\,(0.0015) \\
D2 & 0.0526\,(0.0017) & 0.0534\,(0.0018) & 0.0585\,(0.0019) & \textbf{0.0327\,(0.0007)} & 0.0451\,(0.0012) & 0.0400\,(0.0010) \\
D3 & 0.0703\,(0.0021) & 0.0692\,(0.0019) & 0.0684\,(0.0020) & \textbf{0.0421\,(0.0014)} & 0.0643\,(0.0014) & 0.0597\,(0.0018) \\
D4 & 0.0609\,(0.0020) & 0.0607\,(0.0020) & 0.0599\,(0.0017) & \textbf{0.0233\,(0.0013)} & 0.0405\,(0.0018) & 0.0443\,(0.0027) \\
D5 & 0.0589\,(0.0022) & 0.0592\,(0.0020) & 0.0581\,(0.0016) & \textbf{0.0456\,(0.0010)} & 0.0511\,(0.0013) & 0.0568\,(0.0013) \\
D6 & 0.0839\,(0.0033) & 0.0841\,(0.0034) & 0.0920\,(0.0039) & \textbf{0.0299\,(0.0008)} & 0.0344\,(0.0012) & 0.0400\,(0.0017) \\
D7 & 0.0541\,(0.0019) & 0.0540\,(0.0018) & 0.0590\,(0.0019) & \textbf{0.0345\,(0.0007)} & 0.0398\,(0.0011) & 0.0422\,(0.0011) \\
D8 & \textbf{0.0641\,(0.0018)} & 0.0642\,(0.0018) & 0.0727\,(0.0017) & 0.0972\,(0.0018) & 0.1211\,(0.0030) & 0.1168\,(0.0023) \\
D9 & 0.0735\,(0.0032) & 0.0737\,(0.0028) & 0.0675\,(0.0024) & \textbf{0.0449\,(0.0014)} & 0.0528\,(0.0018) & 0.0658\,(0.0027) \\
\bottomrule
\end{tabular}
\end{table}

\begin{table}[htbp]
\centering
\caption{Stage 2 mode coverage. D6 is the multimodal regime and is the only row that carries information; a uniform 1.0000 elsewhere records that the single-mode regimes stayed single-mode.}
\label{tab:p55s2-mode}
\scriptsize
\setlength{\tabcolsep}{3.2pt}
\begin{tabular}{@{}lrrrrrr@{}}
\toprule
Regime & mutau & R3 & v1 & W-DRF-T & C-DRF & DRF \\
\midrule
D0 & \textbf{1.0000\,(0.0000)} & \textbf{1.0000\,(0.0000)} & \textbf{1.0000\,(0.0000)} & \textbf{1.0000\,(0.0000)} & \textbf{1.0000\,(0.0000)} & \textbf{1.0000\,(0.0000)} \\
D1 & \textbf{1.0000\,(0.0000)} & \textbf{1.0000\,(0.0000)} & \textbf{1.0000\,(0.0000)} & \textbf{1.0000\,(0.0000)} & \textbf{1.0000\,(0.0000)} & \textbf{1.0000\,(0.0000)} \\
D2 & \textbf{1.0000\,(0.0000)} & \textbf{1.0000\,(0.0000)} & \textbf{1.0000\,(0.0000)} & \textbf{1.0000\,(0.0000)} & \textbf{1.0000\,(0.0000)} & \textbf{1.0000\,(0.0000)} \\
D3 & \textbf{1.0000\,(0.0000)} & \textbf{1.0000\,(0.0000)} & \textbf{1.0000\,(0.0000)} & \textbf{1.0000\,(0.0000)} & \textbf{1.0000\,(0.0000)} & \textbf{1.0000\,(0.0000)} \\
D4 & \textbf{1.0000\,(0.0000)} & \textbf{1.0000\,(0.0000)} & \textbf{1.0000\,(0.0000)} & \textbf{1.0000\,(0.0000)} & \textbf{1.0000\,(0.0000)} & \textbf{1.0000\,(0.0000)} \\
D5 & \textbf{1.0000\,(0.0000)} & \textbf{1.0000\,(0.0000)} & \textbf{1.0000\,(0.0000)} & \textbf{1.0000\,(0.0000)} & \textbf{1.0000\,(0.0000)} & \textbf{1.0000\,(0.0000)} \\
D6 & 0.9930\,(0.0012) & 0.9931\,(0.0012) & 0.9897\,(0.0020) & \textbf{1.0000\,(0.0000)} & \textbf{1.0000\,(0.0000)} & \textbf{1.0000\,(0.0000)} \\
D7 & \textbf{1.0000\,(0.0000)} & \textbf{1.0000\,(0.0000)} & \textbf{1.0000\,(0.0000)} & \textbf{1.0000\,(0.0000)} & \textbf{1.0000\,(0.0000)} & \textbf{1.0000\,(0.0000)} \\
D8 & \textbf{1.0000\,(0.0000)} & \textbf{1.0000\,(0.0000)} & \textbf{1.0000\,(0.0000)} & \textbf{1.0000\,(0.0000)} & \textbf{1.0000\,(0.0000)} & \textbf{1.0000\,(0.0000)} \\
D9 & \textbf{1.0000\,(0.0000)} & \textbf{1.0000\,(0.0000)} & \textbf{1.0000\,(0.0000)} & \textbf{1.0000\,(0.0000)} & \textbf{1.0000\,(0.0000)} & \textbf{1.0000\,(0.0000)} \\
\bottomrule
\end{tabular}
\end{table}

\begin{table}[htbp]
\centering
\caption{Stage 2 upper-tail calibration.}
\label{tab:p55s2-tail}
\scriptsize
\setlength{\tabcolsep}{3.2pt}
\begin{tabular}{@{}lrrrrrr@{}}
\toprule
Regime & mutau & R3 & v1 & W-DRF-T & C-DRF & DRF \\
\midrule
D0 & 0.0711\,(0.0020) & \textbf{0.0708\,(0.0022)} & 0.0716\,(0.0019) & 0.1435\,(0.0032) & 0.1783\,(0.0029) & 0.1792\,(0.0041) \\
D1 & 0.1058\,(0.0024) & 0.1058\,(0.0022) & 0.1065\,(0.0024) & \textbf{0.0833\,(0.0016)} & 0.1085\,(0.0025) & 0.1032\,(0.0024) \\
D2 & 0.1135\,(0.0023) & 0.1139\,(0.0024) & 0.1174\,(0.0022) & \textbf{0.0796\,(0.0011)} & 0.1019\,(0.0017) & 0.0938\,(0.0018) \\
D3 & 0.1152\,(0.0022) & 0.1150\,(0.0021) & 0.1134\,(0.0022) & \textbf{0.0877\,(0.0022)} & 0.1259\,(0.0023) & 0.1135\,(0.0029) \\
D4 & 0.0764\,(0.0027) & 0.0771\,(0.0026) & 0.0787\,(0.0020) & \textbf{0.0414\,(0.0018)} & 0.0858\,(0.0032) & 0.0689\,(0.0035) \\
D5 & 0.0953\,(0.0027) & 0.0944\,(0.0025) & 0.0936\,(0.0029) & \textbf{0.0690\,(0.0015)} & 0.0722\,(0.0016) & 0.0805\,(0.0020) \\
D6 & 0.0940\,(0.0033) & 0.0955\,(0.0034) & 0.1010\,(0.0040) & 0.0547\,(0.0013) & 0.0555\,(0.0013) & \textbf{0.0546\,(0.0015)} \\
D7 & 0.0519\,(0.0011) & 0.0521\,(0.0011) & 0.0541\,(0.0011) & \textbf{0.0372\,(0.0010)} & 0.0426\,(0.0010) & 0.0430\,(0.0012) \\
D8 & \textbf{0.0977\,(0.0024)} & 0.0992\,(0.0022) & 0.1027\,(0.0026) & 0.1131\,(0.0022) & 0.1452\,(0.0031) & 0.1320\,(0.0026) \\
D9 & 0.1175\,(0.0036) & 0.1187\,(0.0031) & 0.1153\,(0.0028) & \textbf{0.1050\,(0.0024)} & 0.1124\,(0.0025) & 0.1324\,(0.0031) \\
\bottomrule
\end{tabular}
\end{table}

\begin{table}[htbp]
\centering
\caption{Stage 2 conditional functional targets, \texttt{TCATE-K}. A blank cell is a target the method cannot supply by contract.}
\label{tab:p55s2-tcate}
\scriptsize
\setlength{\tabcolsep}{3.2pt}
\begin{tabular}{@{}llrrrrr@{}}
\toprule
Functional & Regime & mutau & R3 & PTA-S & C-DRF & DRF \\
\midrule
grid\_mean & D0 & 0.0603\,(0.0012) & 0.0569\,(0.0013) & \textbf{0.0113\,(0.0009)} & 0.1153\,(0.0038) & 0.1123\,(0.0061) \\
 & D1 & 0.1251\,(0.0074) & 0.1185\,(0.0075) & \textbf{0.0630\,(0.0042)} & 0.1696\,(0.0071) & 0.1314\,(0.0051) \\
 & D2 & \textbf{0.0226\,(0.0023)} & 0.0240\,(0.0026) & 0.0336\,(0.0031) & 0.0756\,(0.0056) & 0.0491\,(0.0051) \\
 & D3 & 0.1047\,(0.0045) & 0.1016\,(0.0044) & \textbf{0.0802\,(0.0044)} & 0.1381\,(0.0086) & 0.0993\,(0.0060) \\
 & D4 & 0.1367\,(0.0082) & 0.1165\,(0.0064) & \textbf{0.0633\,(0.0042)} & 0.1250\,(0.0064) & 0.1183\,(0.0081) \\
 & D5 & 0.0346\,(0.0031) & 0.0369\,(0.0034) & \textbf{0.0235\,(0.0022)} & 0.0544\,(0.0047) & 0.0608\,(0.0057) \\
 & D6 & 0.2306\,(0.0133) & 0.2239\,(0.0122) & \textbf{0.1896\,(0.0120)} & 0.1980\,(0.0157) & 0.1957\,(0.0141) \\
 & D7 & \textbf{0.0232\,(0.0023)} & 0.0265\,(0.0025) & 0.0319\,(0.0032) & 0.0606\,(0.0053) & 0.0504\,(0.0053) \\
 & D8 & 0.0566\,(0.0032) & \textbf{0.0531\,(0.0030)} & 0.0590\,(0.0042) & 0.2248\,(0.0125) & 0.1250\,(0.0076) \\
 & D9 & 0.2314\,(0.0109) & 0.2380\,(0.0104) & \textbf{0.1329\,(0.0057)} & 0.2430\,(0.0105) & 0.2523\,(0.0114) \\
\midrule
grid\_sd & D0 & 0.0118\,(0.0007) & 0.0104\,(0.0006) & \textbf{0.0011\,(0.0001)} & 0.0180\,(0.0012) & 0.0173\,(0.0013) \\
 & D1 & 0.0289\,(0.0024) & 0.0294\,(0.0023) & \textbf{0.0166\,(0.0014)} & 0.0276\,(0.0017) & 0.0272\,(0.0020) \\
 & D2 & 0.0165\,(0.0015) & 0.0181\,(0.0014) & \textbf{0.0155\,(0.0013)} & 0.0194\,(0.0015) & 0.0205\,(0.0016) \\
 & D3 & 0.0300\,(0.0021) & 0.0296\,(0.0017) & \textbf{0.0171\,(0.0014)} & 0.0284\,(0.0019) & 0.0254\,(0.0016) \\
 & D4 & 0.0471\,(0.0028) & 0.0402\,(0.0025) & \textbf{0.0325\,(0.0023)} & 0.0503\,(0.0024) & 0.0474\,(0.0027) \\
 & D5 & 0.0432\,(0.0034) & 0.0471\,(0.0030) & \textbf{0.0227\,(0.0028)} & 0.0354\,(0.0024) & 0.0314\,(0.0028) \\
 & D6 & 0.0139\,(0.0012) & 0.0139\,(0.0011) & \textbf{0.0119\,(0.0010)} & 0.0174\,(0.0012) & 0.0155\,(0.0013) \\
 & D7 & 0.0194\,(0.0014) & 0.0191\,(0.0015) & \textbf{0.0163\,(0.0012)} & 0.0210\,(0.0016) & 0.0203\,(0.0016) \\
 & D8 & 0.0189\,(0.0014) & \textbf{0.0178\,(0.0015)} & 0.0197\,(0.0015) & 0.0258\,(0.0017) & 0.0252\,(0.0019) \\
 & D9 & 0.0434\,(0.0032) & 0.0407\,(0.0025) & \textbf{0.0175\,(0.0019)} & 0.0572\,(0.0031) & 0.0642\,(0.0032) \\
\midrule
grid\_skewness & D0 & 0.0102\,(0.0014) & 0.0101\,(0.0015) & n/a & 0.0000\,(0.0000) & \textbf{0.0000\,(0.0000)} \\
 & D1 & 0.0000\,(0.0000) & 0.0000\,(0.0000) & n/a & \textbf{0.0000\,(0.0000)} & 0.0000\,(0.0000) \\
 & D2 & 0.0000\,(0.0000) & 0.0000\,(0.0000) & n/a & \textbf{0.0000\,(0.0000)} & 0.0000\,(0.0000) \\
 & D3 & 0.0000\,(0.0000) & 0.0000\,(0.0000) & n/a & 0.0000\,(0.0000) & \textbf{0.0000\,(0.0000)} \\
 & D4 & 0.0000\,(0.0000) & 0.0000\,(0.0000) & n/a & \textbf{0.0000\,(0.0000)} & 0.0000\,(0.0000) \\
 & D5 & 0.0030\,(0.0012) & 0.0034\,(0.0011) & n/a & 0.0000\,(0.0000) & \textbf{0.0000\,(0.0000)} \\
 & D6 & 0.0000\,(0.0000) & 0.0000\,(0.0000) & n/a & 0.0000\,(0.0000) & \textbf{0.0000\,(0.0000)} \\
 & D7 & 0.1188\,(0.0093) & 0.0904\,(0.0067) & n/a & \textbf{0.0500\,(0.0021)} & 0.0567\,(0.0016) \\
 & D8 & 0.0000\,(0.0000) & 0.0000\,(0.0000) & n/a & 0.0000\,(0.0000) & \textbf{0.0000\,(0.0000)} \\
 & D9 & 0.0005\,(0.0003) & 0.0000\,(0.0000) & n/a & 0.0000\,(0.0000) & \textbf{0.0000\,(0.0000)} \\
\midrule
grid\_upper\_tail\_mean & D0 & 0.0607\,(0.0013) & \textbf{0.0570\,(0.0014)} & n/a & 0.1208\,(0.0042) & 0.1158\,(0.0063) \\
 & D1 & 0.1283\,(0.0073) & \textbf{0.1209\,(0.0073)} & n/a & 0.1760\,(0.0078) & 0.1362\,(0.0054) \\
 & D2 & \textbf{0.0264\,(0.0026)} & 0.0291\,(0.0030) & n/a & 0.0782\,(0.0059) & 0.0532\,(0.0054) \\
 & D3 & 0.1091\,(0.0049) & \textbf{0.1049\,(0.0044)} & n/a & 0.1453\,(0.0088) & 0.1057\,(0.0061) \\
 & D4 & 0.1658\,(0.0095) & \textbf{0.1383\,(0.0076)} & n/a & 0.1576\,(0.0074) & 0.1469\,(0.0100) \\
 & D5 & \textbf{0.0508\,(0.0040)} & 0.0552\,(0.0040) & n/a & 0.0611\,(0.0046) & 0.0683\,(0.0057) \\
 & D6 & 0.2291\,(0.0134) & 0.2224\,(0.0123) & n/a & \textbf{0.1969\,(0.0159)} & 0.1969\,(0.0141) \\
 & D7 & \textbf{0.0286\,(0.0027)} & 0.0314\,(0.0024) & n/a & 0.0635\,(0.0054) & 0.0536\,(0.0055) \\
 & D8 & 0.0606\,(0.0034) & \textbf{0.0566\,(0.0031)} & n/a & 0.2319\,(0.0130) & 0.1300\,(0.0075) \\
 & D9 & \textbf{0.2574\,(0.0098)} & 0.2622\,(0.0099) & n/a & 0.2798\,(0.0115) & 0.2935\,(0.0119) \\
\bottomrule
\end{tabular}
\end{table}

\begin{table}[htbp]
\centering
\caption{Stage 2 marginal functional targets, \texttt{TATE-K}.}
\label{tab:p55s2-tate}
\scriptsize
\setlength{\tabcolsep}{3.2pt}
\begin{tabular}{@{}llrrrrr@{}}
\toprule
Functional & Regime & mutau & R3 & PTA-S & C-DRF & DRF \\
\midrule
grid\_mean & D0 & 0.0057\,(0.0007) & 0.0065\,(0.0009) & \textbf{0.0038\,(0.0005)} & 0.0547\,(0.0039) & 0.0286\,(0.0032) \\
 & D1 & \textbf{0.0169\,(0.0020)} & 0.0181\,(0.0024) & 0.0207\,(0.0021) & 0.0925\,(0.0068) & 0.0319\,(0.0042) \\
 & D2 & \textbf{0.0150\,(0.0015)} & 0.0156\,(0.0019) & 0.0241\,(0.0025) & 0.0686\,(0.0058) & 0.0226\,(0.0042) \\
 & D3 & 0.0216\,(0.0023) & \textbf{0.0177\,(0.0021)} & 0.0228\,(0.0020) & 0.0289\,(0.0043) & 0.0234\,(0.0025) \\
 & D4 & \textbf{0.0187\,(0.0024)} & 0.0196\,(0.0023) & 0.0217\,(0.0023) & 0.0613\,(0.0062) & 0.0312\,(0.0048) \\
 & D5 & 0.0168\,(0.0021) & \textbf{0.0155\,(0.0023)} & 0.0156\,(0.0018) & 0.0384\,(0.0053) & 0.0207\,(0.0036) \\
 & D6 & \textbf{0.0427\,(0.0053)} & 0.0476\,(0.0057) & 0.0956\,(0.0114) & 0.1205\,(0.0147) & 0.1158\,(0.0146) \\
 & D7 & \textbf{0.0145\,(0.0019)} & 0.0152\,(0.0018) & 0.0222\,(0.0024) & 0.0478\,(0.0057) & 0.0234\,(0.0042) \\
 & D8 & 0.0315\,(0.0034) & \textbf{0.0313\,(0.0033)} & 0.0336\,(0.0046) & 0.2001\,(0.0144) & 0.0751\,(0.0086) \\
 & D9 & 0.0301\,(0.0028) & \textbf{0.0247\,(0.0033)} & 0.0353\,(0.0035) & 0.1872\,(0.0106) & 0.2076\,(0.0117) \\
\midrule
grid\_sd & D0 & 0.0064\,(0.0007) & 0.0055\,(0.0007) & \textbf{0.0007\,(0.0001)} & 0.0134\,(0.0012) & 0.0122\,(0.0012) \\
 & D1 & 0.0167\,(0.0019) & 0.0170\,(0.0020) & \textbf{0.0134\,(0.0015)} & 0.0192\,(0.0021) & 0.0181\,(0.0023) \\
 & D2 & \textbf{0.0128\,(0.0015)} & 0.0130\,(0.0012) & 0.0131\,(0.0014) & 0.0140\,(0.0017) & 0.0142\,(0.0016) \\
 & D3 & 0.0157\,(0.0017) & 0.0145\,(0.0016) & \textbf{0.0133\,(0.0015)} & 0.0199\,(0.0020) & 0.0181\,(0.0021) \\
 & D4 & 0.0155\,(0.0018) & \textbf{0.0151\,(0.0018)} & 0.0151\,(0.0014) & 0.0305\,(0.0026) & 0.0223\,(0.0025) \\
 & D5 & 0.0279\,(0.0034) & 0.0260\,(0.0032) & \textbf{0.0219\,(0.0029)} & 0.0247\,(0.0026) & 0.0242\,(0.0027) \\
 & D6 & \textbf{0.0092\,(0.0011)} & 0.0098\,(0.0010) & 0.0093\,(0.0010) & 0.0110\,(0.0014) & 0.0113\,(0.0014) \\
 & D7 & 0.0142\,(0.0012) & 0.0136\,(0.0013) & \textbf{0.0125\,(0.0014)} & 0.0141\,(0.0017) & 0.0141\,(0.0016) \\
 & D8 & 0.0136\,(0.0013) & \textbf{0.0126\,(0.0015)} & 0.0174\,(0.0017) & 0.0186\,(0.0020) & 0.0194\,(0.0022) \\
 & D9 & \textbf{0.0146\,(0.0019)} & 0.0162\,(0.0022) & 0.0150\,(0.0021) & 0.0470\,(0.0037) & 0.0519\,(0.0036) \\
\midrule
grid\_skewness & D0 & 0.0075\,(0.0014) & 0.0085\,(0.0015) & n/a & 0.0000\,(0.0000) & \textbf{0.0000\,(0.0000)} \\
 & D1 & 0.0000\,(0.0000) & 0.0000\,(0.0000) & n/a & \textbf{0.0000\,(0.0000)} & 0.0000\,(0.0000) \\
 & D2 & 0.0000\,(0.0000) & 0.0000\,(0.0000) & n/a & 0.0000\,(0.0000) & \textbf{0.0000\,(0.0000)} \\
 & D3 & 0.0000\,(0.0000) & 0.0000\,(0.0000) & n/a & \textbf{0.0000\,(0.0000)} & 0.0000\,(0.0000) \\
 & D4 & 0.0000\,(0.0000) & 0.0000\,(0.0000) & n/a & \textbf{0.0000\,(0.0000)} & 0.0000\,(0.0000) \\
 & D5 & 0.0023\,(0.0011) & 0.0029\,(0.0011) & n/a & 0.0000\,(0.0000) & \textbf{0.0000\,(0.0000)} \\
 & D6 & 0.0000\,(0.0000) & 0.0000\,(0.0000) & n/a & \textbf{0.0000\,(0.0000)} & 0.0000\,(0.0000) \\
 & D7 & 0.1024\,(0.0084) & 0.0746\,(0.0060) & n/a & 0.0167\,(0.0016) & \textbf{0.0130\,(0.0014)} \\
 & D8 & 0.0000\,(0.0000) & 0.0000\,(0.0000) & n/a & 0.0000\,(0.0000) & \textbf{0.0000\,(0.0000)} \\
 & D9 & 0.0005\,(0.0003) & 0.0000\,(0.0000) & n/a & 0.0000\,(0.0000) & \textbf{0.0000\,(0.0000)} \\
\midrule
grid\_upper\_tail\_mean & D0 & \textbf{0.0070\,(0.0009)} & 0.0078\,(0.0010) & n/a & 0.0644\,(0.0041) & 0.0364\,(0.0035) \\
 & D1 & \textbf{0.0223\,(0.0027)} & 0.0234\,(0.0031) & n/a & 0.1037\,(0.0077) & 0.0410\,(0.0052) \\
 & D2 & \textbf{0.0188\,(0.0021)} & 0.0197\,(0.0024) & n/a & 0.0707\,(0.0062) & 0.0261\,(0.0046) \\
 & D3 & 0.0247\,(0.0026) & \textbf{0.0223\,(0.0023)} & n/a & 0.0389\,(0.0054) & 0.0288\,(0.0033) \\
 & D4 & 0.0257\,(0.0029) & \textbf{0.0244\,(0.0029)} & n/a & 0.0837\,(0.0071) & 0.0440\,(0.0058) \\
 & D5 & 0.0317\,(0.0029) & 0.0307\,(0.0029) & n/a & 0.0426\,(0.0057) & \textbf{0.0299\,(0.0039)} \\
 & D6 & \textbf{0.0424\,(0.0049)} & 0.0477\,(0.0056) & n/a & 0.1189\,(0.0149) & 0.1145\,(0.0148) \\
 & D7 & \textbf{0.0194\,(0.0025)} & 0.0195\,(0.0023) & n/a & 0.0503\,(0.0058) & 0.0260\,(0.0046) \\
 & D8 & \textbf{0.0328\,(0.0038)} & 0.0339\,(0.0034) & n/a & 0.2076\,(0.0149) & 0.0833\,(0.0088) \\
 & D9 & 0.0310\,(0.0032) & \textbf{0.0282\,(0.0042)} & n/a & 0.2233\,(0.0109) & 0.2479\,(0.0121) \\
\bottomrule
\end{tabular}
\end{table}

\begin{table}[htbp]
\centering
\caption{Stage 2 seed-paired differences on the primary law metric, \texttt{kernel\_law\_error}.}
\label{tab:p55s2-paired-law}
\scriptsize
\setlength{\tabcolsep}{4pt}
\begin{tabular}{@{}lrrrrrrrrr@{}}
\toprule
Regime & \multicolumn{3}{c}{vs R3} & \multicolumn{3}{c}{vs v1} & \multicolumn{3}{c}{vs C-DRF} \\
  & diff. & SE & $|$diff$|/$SE & diff. & SE & $|$diff$|/$SE & diff. & SE & $|$diff$|/$SE \\
\midrule
D0 & +0.0004 & 0.0005 & 0.7 & +0.0032$^{\dagger}$ & 0.0007 & 4.5 & -0.0883$^{\ast}$ & 0.0027 & 33.0 \\
D1 & +0.0002 & 0.0001 & 1.4 & +0.0006$^{\dagger}$ & 0.0002 & 2.8 & -0.0052$^{\ast}$ & 0.0007 & 7.8 \\
D2 & -0.0002 & 0.0002 & 0.8 & -0.0013$^{\ast}$ & 0.0002 & 6.2 & -0.0070$^{\ast}$ & 0.0004 & 18.7 \\
D3 & +0.0005$^{\dagger}$ & 0.0002 & 3.1 & +0.0003 & 0.0002 & 1.3 & -0.0091$^{\ast}$ & 0.0003 & 28.4 \\
D4 & +0.0003$^{\dagger}$ & 0.0001 & 2.2 & +0.0007$^{\dagger}$ & 0.0002 & 4.2 & -0.0010$^{\ast}$ & 0.0004 & 2.5 \\
D5 & +0.0016$^{\dagger}$ & 0.0007 & 2.4 & +0.0016$^{\dagger}$ & 0.0006 & 2.8 & -0.0180$^{\ast}$ & 0.0009 & 19.6 \\
D6 & -0.0002 & 0.0002 & 1.1 & -0.0034$^{\ast}$ & 0.0005 & 6.9 & +0.0167$^{\dagger}$ & 0.0010 & 16.1 \\
D7 & +0.0001 & 0.0002 & 0.9 & -0.0012$^{\ast}$ & 0.0002 & 6.2 & -0.0040$^{\ast}$ & 0.0003 & 13.1 \\
D8 & -0.0000 & 0.0002 & 0.0 & -0.0025$^{\ast}$ & 0.0002 & 11.1 & -0.0337$^{\ast}$ & 0.0011 & 30.6 \\
D9 & +0.0002 & 0.0006 & 0.3 & +0.0031$^{\dagger}$ & 0.0005 & 5.8 & -0.0001 & 0.0006 & 0.1 \\
\bottomrule
\end{tabular}
\par\vspace{2pt}\parbox{0.96\textwidth}{\scriptsize Negative favours \texttt{mutau}. $\ast$ marks a claimant win and $\dagger$ a comparator win, both at more than two paired standard errors, the frozen decision multiple. The third column of each block expresses the same difference in standard errors, so a reader can see how far past the threshold a result sits rather than only that it crossed.}
\end{table}

\begin{table}[htbp]
\centering
\caption{Stage 2 seed-paired differences on the grid causal mean, \texttt{mean\_quantile\_rmse}.}
\label{tab:p55s2-paired-meanq}
\scriptsize
\setlength{\tabcolsep}{4pt}
\begin{tabular}{@{}lrrrrrrrrr@{}}
\toprule
Regime & \multicolumn{3}{c}{vs R3} & \multicolumn{3}{c}{vs PTA-S} & \multicolumn{3}{c}{vs C-DRF} \\
  & diff. & SE & $|$diff$|/$SE & diff. & SE & $|$diff$|/$SE & diff. & SE & $|$diff$|/$SE \\
\midrule
D0 & +0.0025$^{\dagger}$ & 0.0009 & 2.7 & +0.0430$^{\dagger}$ & 0.0008 & 55.9 & -0.0643$^{\ast}$ & 0.0028 & 23.2 \\
D1 & +0.0009 & 0.0015 & 0.6 & +0.0626$^{\dagger}$ & 0.0031 & 20.0 & -0.0066 & 0.0040 & 1.7 \\
D2 & -0.0073 & 0.0080 & 0.9 & -0.0112$^{\ast}$ & 0.0049 & 2.3 & -0.0527$^{\ast}$ & 0.0067 & 7.8 \\
D3 & +0.0035$^{\dagger}$ & 0.0015 & 2.3 & -0.0050$^{\ast}$ & 0.0022 & 2.3 & -0.1405$^{\ast}$ & 0.0027 & 51.2 \\
D4 & +0.0060$^{\dagger}$ & 0.0017 & 3.6 & +0.0690$^{\dagger}$ & 0.0037 & 18.7 & +0.0459$^{\dagger}$ & 0.0033 & 13.8 \\
D5 & -0.0298$^{\ast}$ & 0.0097 & 3.1 & +0.0846$^{\dagger}$ & 0.0121 & 7.0 & +0.0461$^{\dagger}$ & 0.0122 & 3.8 \\
D6 & +0.0053 & 0.0107 & 0.5 & +0.0210$^{\dagger}$ & 0.0094 & 2.2 & -0.0036 & 0.0099 & 0.4 \\
D7 & +0.0030 & 0.0043 & 0.7 & +0.0077 & 0.0051 & 1.5 & -0.0157$^{\ast}$ & 0.0058 & 2.7 \\
D8 & +0.0042$^{\dagger}$ & 0.0016 & 2.7 & -0.0173$^{\ast}$ & 0.0028 & 6.1 & -0.1861$^{\ast}$ & 0.0105 & 17.8 \\
D9 & -0.0000 & 0.0079 & 0.0 & +0.0919$^{\dagger}$ & 0.0066 & 14.0 & -0.0072 & 0.0058 & 1.2 \\
\bottomrule
\end{tabular}
\par\vspace{2pt}\parbox{0.96\textwidth}{\scriptsize Negative favours \texttt{mutau}. $\ast$ marks a claimant win and $\dagger$ a comparator win, both at more than two paired standard errors, the frozen decision multiple. The third column of each block expresses the same difference in standard errors, so a reader can see how far past the threshold a result sits rather than only that it crossed.}
\end{table}

\begin{table}[htbp]
\centering
\caption{The degradation clause, per sample size. Paired \texttt{kernel\_law\_error} difference, \texttt{mutau} minus C-WDB R3, on the two regimes the phase document names.}
\label{tab:p55s2-degradation}
\scriptsize
\begin{tabular}{@{}llrrrl@{}}
\toprule
Regime & $n$ & diff. & SE & $|$diff$|/$SE & Verdict \\
\midrule
D3 & 500 & +0.00078 & 0.00026 & 3.0 & \textbf{R3 better} \\
D3 & 1000 & +0.00026 & 0.00020 & 1.3 & tie \\
D7 & 500 & +0.00030 & 0.00026 & 1.2 & tie \\
D7 & 1000 & -0.00004 & 0.00016 & 0.2 & tie \\
\bottomrule
\end{tabular}
\par\vspace{2pt}\parbox{0.9\textwidth}{\scriptsize D3 is the separate-head-favourable regime, which is what the $\mu/\tau$ reparameterisation puts at risk because it carries the prognostic and contrast fields in one shared tree. D7 is the pure-shape-transfer regime. Positive favours R3.}
\end{table}

\begin{table}[htbp]
\centering
\caption{Mechanism diagnostics for \texttt{cwdb\_mutau} across all ten regimes, pooled over sample sizes and twenty seeds.}
\label{tab:p55s2-diagnostics}
\scriptsize
\begin{tabular}{@{}lrrrr@{}}
\toprule
Regime & Effective support (of $M = 10$) & Selected contrast strength & Boosting steps & Training risk \\
\midrule
D0 & 10.00 & 0.0 & 100 & 0.0227 \\
D1 & 10.00 & 28.8 & 100 & 0.1233 \\
D2 & 10.00 & 341.2 & 100 & 0.1314 \\
D3 & 10.00 & 0.0 & 100 & 0.1097 \\
D4 & 10.00 & 33.8 & 100 & 0.1335 \\
D5 & 10.00 & 32.5 & 100 & 0.1339 \\
D6 & 10.00 & 185.0 & 100 & 0.6427 \\
D7 & 10.00 & 295.0 & 100 & 0.1307 \\
D8 & 10.00 & 320.0 & 100 & 0.1587 \\
D9 & 10.00 & 102.5 & 100 & 0.1503 \\
\bottomrule
\end{tabular}
\par\vspace{2pt}\parbox{0.9\textwidth}{\scriptsize The rule 5 floor is $0.6M = 6$ particles. The selected contrast strength is chosen per cell by held-out risk over the frozen candidate set $\{0, 50, 500\}$, so a regime mean between the candidates means the selection varies across cells rather than that an intermediate value was ever used.}
\end{table}

\begin{table}[htbp]
\centering
\caption{Runtime on the main grid. The \texttt{mutau} row is measured on the Stage 2 rows alone; every other row is that method's own frozen record.}
\label{tab:p55s2-cost}
\scriptsize
\begin{tabular}{@{}lrrr@{}}
\toprule
Method & Median (s) & Max (s) & Ratio to Causal-DRF \\
\midrule
DRF & 11.90 & 41.1 & 0.8 \\
C-DRF & 14.36 & 46.2 & 1.0 \\
v1 & 16.18 & 36.3 & 1.1 \\
PTA-S & 36.20 & 52.8 & 2.5 \\
R3 & 92.23 & 135.7 & 6.4 \\
mutau & 199.27 & 420.3 & 13.9 \\
\bottomrule
\end{tabular}
\par\vspace{2pt}\parbox{0.92\textwidth}{\scriptsize The rule 6 cap is 60. The denominator is the original-code Causal-DRF median of 14.36\,s, declared in the preregistration before the run because the project contains two Causal-DRF records and the verdict depends on which is used. Every figure here is an upper bound: the Stage 2 cells ran sixteen at a time on one laptop, and the same implementation measured on an idle machine completes a cell in about a fifth of the time. Runtime in this project is dominated by machine load and is not a clean measure of algorithmic cost.}
\end{table}

